\documentclass[11pt]{article}
% Diplomatic transcription macros (12 Nov 2026). Typed text is plain; the handwritten layer is blue,
% deletions red, typewriter strike-outs grey, doubtful readings in brackets with a question mark.
\usepackage[utf8]{inputenc}
\usepackage[T1]{fontenc}
\usepackage{lmodern}
\usepackage[a4paper,margin=2.2cm]{geometry}
\usepackage{amsmath,amssymb}
\usepackage{xcolor}
\usepackage{tikz-cd}
\usepackage{environ}
\setlength{\parindent}{0pt}
\setlength{\parskip}{4pt}
\title{Tapis de Quillen, with Quillen's offprint and preprint (cote 162-5, typed part)\\[4pt]\large Diplomatic transcription from the Montpellier scans}
\author{AI-transcribed by Claude Opus 5.5}
\date{12 November 2026}
\newcommand{\scanpage}[1]{\par\bigskip\noindent{\small\sffamily[scan #1]}\par}
\newcommand{\mspage}[1]{{\small\sffamily[p.\,#1]}}
\newcommand{\hand}[1]{\textcolor{blue}{#1}}
\let\over\relax
\newcommand{\over}[2]{\textcolor{gray}{#1}\textcolor{blue}{\ensuremath{^{\!\nearrow}}#2}}
\newcommand{\inserted}[2][]{\textcolor{blue}{$\langle$#2$\rangle$}}
\newcommand{\margin}[1]{\textcolor{blue}{\{\textsf{marg.:} #1\}}}
\newcommand{\deleted}[1]{\textcolor{red}{[\textsf{del.:} #1]}}
\newcommand{\xout}[1]{\textcolor{gray}{[\textsf{x:} #1]}}
\newcommand{\uline}[1]{\underline{#1}}
\newcommand{\uuline}[1]{\underline{\underline{#1}}}
\newcommand{\unsure}[1]{[#1?]}
\newcommand{\illegible}[1]{[\textsf{illeg.\ #1}]}
\newcommand{\diagram}[1]{\par{\centering #1\par}}
\newenvironment{diagramblock}{\par\begin{center}}{\end{center}\par}

\begin{document}
\maketitle

% Tapis de Quillen, with Quillen's offprint and preprint (cote 162-5, typed part): archive cote 162-5, site page 2 (PDF page 3).
% typescript. Draft compared: none.
% AI-transcribed by Claude Opus 5.5 (claude-opus-5-5, adaptive thinking, effort high), Message Batches API, batch t162_full (msgbatch_01UAh3zyWPd454AeDkhnmQ33), 2026-11-13; not yet checked by eye.
\scanpage{162-5/2}
\mspage{\hand{1}} \margin{notes du \uline{10.9.68}}

Tapis de Quillen.

I \quad \uline{Relations entre catégories et ensembles semi-simplic\unsure{ia}ux.}

A toute catégorie C, \ on associe un ensemble semi-simplicial S(C),
trouvant ainsi un foncteur pleinement fidèle \xout{de Cat}
\[ S;\ (\mathrm{Cat}) \longrightarrow (\mathrm{Ssimpl}) \quad . \]
\xout{Si C est associé à} Les systêmes locaux d'ensemble sur SC correspondent
aux foncteurs sur C qui transforment toute flêche en isomorphisme (i.e.
qu\unsure{i} se factorisent par le groupoïde associé à \ C). Les H$^{i}$ sur SC d'un tel
système local (H$^{o}$ pour ens, H$^{1}$ pour groupes, H$^{i}$ quelconques pour groupes
abéliens) s'interprètent en termes des foncteurs $\uline{\mathrm{lim}}^{(i)}$ dérigés de $\uline{\mathrm{lim}}$ ,
ou si on préfère, des H$^{i}$( du \uline{topos} C\hand{\unsure{$\widetilde{\ }$}}. On voit ainsi à quelle condition
un foncteur \ C $\longrightarrow$ C' induit un\xout{e} \ homotopisme \ SC $\longrightarrow$ SC' : \ en vertu du
critère cohomologique de Artin-Mazur, il f et s que \ pour tout systême
de coéfficients \ F'sur C', l'homomorphisme naturel $\underset{C'}{\underleftarrow{\mathrm{lim}}}^{(i)}F' \longrightarrow \underset{C}{\underleftarrow{\mathrm{lim}}}^{(i\ )}F$
soit un isomorphisme (pour les i pour lesquels cela a un sens). \xout{\unsure{Exemples}}

A C on peut asso\unsure{ci}er le topos $\hand{\widetilde{C}}$ , qui varie de façon \uline{covariante} avec C
% NOTE: right edge of the scan cuts off the ends of this and the following lines; the visible letters are transcribed
(NB le foncteur C$\longmapsto$$\hand{\widetilde{C}}$ \ n'a plus rien de pleinement fidèle, semble-t\unsure{-}il ??)
Les systêmes de coéfficients \inserted{ensemblistes} sur C ( \inserted{déf} = \ les foncteurs \inserted{$C^{o} \longrightarrow$ Ens} transformant isomo\unsure{r-}%
phismes en isomorphismes) correspondent aux faisceaux localement constan\unsure{ts}
i.e. les objets localement constants de $\hand{\widetilde{C}}$, définis intrinsèquement en te\unsure{rmes}
de $\hand{\widetilde{C}}$ . \ Ainsi, le fait pour un foncteur F:C$\longrightarrow$C' d'induire un\xout{e équivalence}
homotopisme S(C)$\longrightarrow$S(C') \ ne \ dépend que du morphisme de topos $\hand{\widetilde{F}}$: $\hand{\widetilde{C}} \longrightarrow \hand{\widetilde{C'}}$
induit, et signifie que pour tout faisceau localement constant F' sur C'
c. sur $\hand{\widetilde{C'}}$ , \ les applic\unsure{a}tions induites \ H$^{i}$($\hand{\widetilde{C'}}$,F') $\longrightarrow$ H$^{i}$($\hand{\widetilde{C}}$,$\hand{\widetilde{F}}^{*}$(F'))
sont des isomorphismes (pour les i pour lesquels cela a un\xout{\unsure{es}} sens).

\margin{2}

\clearpage
% Tapis de Quillen, with Quillen's offprint and preprint (cote 162-5, typed part): archive cote 162-5, site page 3 (PDF page 4).
% typescript. Draft compared: none.
% AI-transcribed by Claude Opus 5.5 (claude-opus-5-5, adaptive thinking, effort high), Message Batches API, batch t162_full (msgbatch_01UAh3zyWPd454AeDkhnmQ33), 2026-11-13; not yet checked by eye.
\scanpage{162-5/3}
\mspage{\hand{2}}

On a aussi un foncteur évident
\[ T: (\mathrm{Ssimpl}) \longrightarrow (\mathrm{Cat}) , \]
en associant à tout ensemble semi-simplicial X la catégorie $T(X) = \hand{\Delta}_{/X}$
des simplexes sur X, dont l'ensemble des objets est la réunion d\unsure{i|a}sjointe
des $X_n$ .... (c'est une catégorie fibrée sur la catégorie $\hand{\Delta}$ des simplexes
types, à fibres les catégories discrètes définies par les $X_n$). Ceci posé
Quillen prouve que pour tout X, ST(X) est isomorphe canoniquement à X da
% NOTE: the typed line runs off the right edge of the sheet after "da"
la catégorie homotopique construite avec (Ssimpl), et que pour toute C,
la catégorie TS(C) est \inserted[typed]{canoniquement} ``homotopiquement équivalente à C'' \xout{\unsure{dansxlaxxxxx}} i.e
\inserted[typed]{\unsure{canoniquement}} isomorphe à C dans la catégorie quotient de (Cat) obtenue en inversant le
% NOTE: the typed line runs off the right edge of the sheet after "le"
foncteurs qui sont des homotopismes. \xout{Ilxenxrésulte} Ces isomorphismes son
% NOTE: the typed line runs off the right edge of the sheet after "son"
fonctoriels en X. Il en résulte formellement qu'un morphisme $f:X\rightarrow Y$ dans
(Ssimpl) est un homotopisme si et seulement si il en est ainsi de $T(f):T($
% NOTE: the typed line runs off the right edge of the sheet after "T("
$\rightarrow T(Y)$, d'où des foncteurs $S'$: $(\mathrm{Cat})' \longrightarrow (\mathrm{Ssimpl})'$ et $T'$:$(\mathrm{Ssimpl})' \longrightarrow (\mathrm{C}$
% NOTE: the typed line runs off the right edge of the sheet after "(C"
entre les catégories ``homotopiques'', \inserted[typed]{construites avec (Cat) resp. (Ssimpl),} qui sont quasi-inverses l'un de l'au
re.

De plus, Quillen construit un isomorphisme canonique et fonctoriel
dans (Cat)' entre C et la catégorie opposée $C^{o}$ , ou ce qui revient au
même, un isomorphisme canonique et fonctoriel dans (Ssimpl)' entre S(C) e
% NOTE: the typed line runs off the right edge of the sheet after "e"
$S(C^{o})$. \xout{\unsure{Celaxluixpermetxalors,xpourx}} La définition est telle que le
foncteur induit sur les systèmes locaux sur C transforme le foncteur
contravariant F sur C, transformant toute flêche en flêche inversible, en
le foncteur covariant (i.e. contravariant sur $C^{o}$) ayant mêmes valeurs sur
les objets, et obtenu sur les flêches en remplaçant $F(u)$ par $F(u)^{-1}$; en
d'autres termes, l'effet de l'homotopisme de Quillen sur les groupoides
fondamentaux est l'isomorphisme évident entre les groupoides fondamentaux
de C et de $C^{o}$ ,\xout{\unsure{quixestxlixdxxx}} compte tenu que le deuxième est l'opposé d
% NOTE: the typed line runs off the right edge of the sheet after "d"

\margin{\hand{3}}

\clearpage
% Tapis de Quillen, with Quillen's offprint and preprint (cote 162-5, typed part): archive cote 162-5, site page 4 (PDF page 5).
% typescript. Draft compared: none.
% AI-transcribed by Claude Opus 5.5 (claude-opus-5-5, adaptive thinking, effort high), Message Batches API, batch t162_full (msgbatch_01UAh3zyWPd454AeDkhnmQ33), 2026-11-13; not yet checked by eye.
\scanpage{162-5/4}
\mspage{\hand{3}}

% NOTE: the right edge of the page is cut off on the scan; the last letters of most lines are not visible
premier. Comme application, Quillen obtient une interprétation faisceau-%
que de la cohomologie d'un ensemble semi-simplicial à coéfficients dans
système \inserted{local} \unsure{contravariant} \inserted{F} \xout{\unsure{de}}(défini classiquement par le complexe cosimplici-
\margin{$\Pi$ !} des $C^n(F) = \coprod_{x \in X_n} F(x)$ ): on considère le système local contravariant
défini par F, on l'interprète comme un faisceau sur \inserted[typed]{T(X) i.e. objet de} $(\mathrm{Ssimpl})_{/X}$ , et on
prend sa cohomologie. - Cependant, quand F est un système de coéfficients
\margin{[ (bracket in the left margin, from this line to ``faisceautique en termes classiques.'')} covariant pas nécessairement local, on n'a tjrs pas d'interprétation de
\margin{Si !} groupes de cohomologie classiques en t ermes \unsure{faiscdautiques|faisceautiques}; ni, lorsque
est \unsure{contravariant}, de son homologie, ou inversement de sa cohomologie
faisceautique en termes classiques.

A propos de la notion de foncteur qui est un homotopisme. Quillen
\margin{| (vertical line in the left margin, from this line to ``logue pour C'')} montre qu\xout{e}'un tel foncteur $F\colon C \to C'$ induit une équivalence entre la sous-
catégorie triangulée \inserted{$D^b_{lc}(C')$} de la catégorie \inserted{dérivée \inserted{bornée} de celle \deleted{des}} des faisceaux abéliens sur C', dont
les faisceaux de cohomologie sont des systèmes locaux, et la catégorie ana-
logue pour C; \xout{\unsure{Il existe}} et réciproquement. \xout{Pour la partie directe} On peut
dans cet énoncé introduire aussi n'importe quel anneau de base (à condition
de le supposer $\neq 0$ dans le cas de la réciproque); \xout{constant} la partie directe
vaut aussi avec un anneau de boefficients pas néc. constant, mais constant
tordu. Je pense que ce résultat (facile) doit pouvoir se généraliser ainsi
Soit $f\colon X \quad X'$ un morphisme de topos qui soit tel que pour tout \xout{systè}
faisceau loc. const. sur X', f induise un isomorphisme sur les cohomologies
(avec cas non commutatif inclus). Supposons que \deleted{\illegible{1} (X)} \inserted{\unsure{X} et X'} soit \uline{locale}
\margin{\o{} (call sign)} \margin{|| (double vertical line in the left margin, from this line to ``nuls pour'')} \uline{homotopiquement trivial}, i.e. que \xout{tout} pour tout entier n \inserted{$\geq 1$}, tout objet U de
\hand{\unsure{X'}} ait un \xout{\illegible{1}} \inserted[typed]{recouvrement} par des $U_i \to U$, tels que a) tout système local sur U
devient constant sur $U_i$ , et toute section sur U devient constante sur $U_i$
et b) pour tout \xout{système} groupe abélien G, les $H^j(U,G) \longrightarrow H^j(U_i,G)$ sont
nuls pour $1 \leq j \leq n$. Alors le foncteur $D^b_{lc}(X') \longrightarrow D^b_{lc}(X)$ induit par f est

\margin{\o{}} \hand{Attention, cette condition n'est \unsure{pratiquement} \unsure{pas} satisfaite par les schémas, avec leur topologie étale (mais bien par \deleted{\illegible{1}} \unsure{(noeth.} avec top. Zariski).}

\margin{4}

\clearpage
% Tapis de Quillen, with Quillen's offprint and preprint (cote 162-5, typed part): archive cote 162-5, site page 5 (PDF page 6).
% typescript. Draft compared: none.
% AI-transcribed by Claude Opus 5.5 (claude-opus-5-5, adaptive thinking, effort high), Message Batches API, batch t162_full (msgbatch_01UAh3zyWPd454AeDkhnmQ33), 2026-11-13; not yet checked by eye.
\scanpage{162-5/5}
\mspage{\hand{4}}

\xout{pleinementxfidèlex(respx} une équivalence). Même énoncé si on met dans le
coup un \xout{Anneauxlocx} système local d'anneaux sur X'. Enfin, f induit \xout{unx}
\xout{fonctepxpleinxfidxx(respx} une équivalence).
\xout{\unsure{isomorphismx}} entre la catégorie des coefficients locaux sur C et celle
\inserted[typed]{\xout{cettexpropriétéxdoit}}
des coéff. locaux sur C' \xout{(??xpeutxêtrexxlxxdoitxilxêtrexmisxdansxlesx}
\xout{hypothèsesxpréliminairesxsurxfx?)x} Principe de démonstration: on commence
par prouver ce dernier résultat, en notant que \xout{dansxlexcax} si un topos est
localement hom. trivial, il est loc\unsure{a} connexe et loc. simplement connexe,
d'où une bonne théorie du $\pi_0$ et du $\pi_1$ (qui sont ici discrets), et on est
ramené \xout{pour} à un cas particulier du critère d'homotopisme de Artin et
Mazur, savoir un critère \inserted[typed]{cohomologique} \unsure{(}pour qu'un homomorhisme de groupes $G\to H$ (ici de
groupes $\pi_1$ de X,X') soit un isomorphisme: il doit induire des isomorphis-%
mes sur les $H^{o}$ et $H^{1}$ (y inclus dans le cas non commutatif...). On prouve
la pleine fidélité en se ramenant par la suite spectrale aux cas des $\mathrm{Ext}^{\unsure{i}}$
de coéfficients locaux. Par la suite spectrale encore, cela résultera du
fait suivant: si X\deleted{'} est \unsure{localement} hom. trivial, alors la catégorie des
faisceaux abéliens loc. constants  est stable par $\underline{\mathrm{Ext}}^{i}$ , et le foncteur
$M\mapsto M_{X}$ \xout{desxfaisxxxxxxxconstants} (Ab) dans $X_{ab}$ commute aux dits $\mathrm{Ext}^{i}$ .
En fait, X et X' étant loc.hom\unsure{p}\hand{|}triviaux, les conditions suivantes sur f
seront équivalentes:

a) f est\xout{'} un homotopisme, i.e.  induit pour tout système local (pas
néc commutatif) sur X'un isomorphisme sur les $H^{i}$\xout{pourxixxm}.

b) \xout{Mêmexchosexave} f induit une équivalence $D^{b}_{lc}(X')\longrightarrow D^{b}_{lc}(X)$ .

a\hand{'}) f induit une équivalence entre la catégorie des systêmes locaux
abéliens sur X' et X, et des isomorphismes sur les $H^{i}$ correspondants (do
on ne prend ici que des coéfficients commutatifs).

\margin{?} J'ignore si on peut dans a) se borner aux systêmes locaux commutatif

\xout{ixexxdanx}
L'équivalence entre a) et b) fournit une première justification ou
motivation pour définir des types d'homotopie via la catégorie D\hand{$^{b}_{lc}$(X)},

\margin{5}

\clearpage
% Tapis de Quillen, with Quillen's offprint and preprint (cote 162-5, typed part): archive cote 162-5, site page 6 (PDF page 7).
% typescript. Draft compared: none.
% AI-transcribed by Claude Opus 5.5 (claude-opus-5-5, adaptive thinking, effort high), Message Batches API, batch t162_full (msgbatch_01UAh3zyWPd454AeDkhnmQ33), 2026-11-13; not yet checked by eye.
\scanpage{162-5/6}
\mspage{5}

éventuellement muni de la sous-catégorie pleine de tous les systêmes loc
% NOTE: the right edge of the scan cuts off the end of this and several following lines
sur X, et du foncteur cohomologique sur $D^{b}_{lc}(X)$ à valeurs \unsure{dqns|dans} \unsure{laddtë|ladite}
catégorie, et bien sûr du produit tensoriel (mais alors on sort de $D^{b}$
pour entrer dans $D^{-}$ , redactor demerdetur).

II \quad n-catégories, catégories n-uples, et (Gr)-catégories.

Appelons bi-catégorie une catégorie dans(Cat), i.e. un triple
($C^{\deleted{o}}_{\hand{0}}$,$C^{\deleted{1}}_{\hand{1}}$, \xout{\unsure{xxoïp}}$p_{1}$,$p_{2}$, $\pi$ ), où $C^{o}$ est une catégorie, $p_{1}$,$p_{2}$: $C^{\deleted{1}}_{\hand{1}} \rightrightarrows C^{\deleted{o}}_{\hand{0}}$ deux
foncteurs, $\pi$:$(C^{\deleted{1}}_{\hand{1}},p_{1})\mathrm{x}_{C^{o}}(C^{\deleted{1}}_{\hand{1}},p_{2}) \longrightarrow C^{\deleted{1}}_{\hand{1}}$ un foncteur, avec les axiomes
habituels. Interprétant (Cat) comme la sous-catégorie pleine de (Ssimpl
des objets qui .... , et de même pour les catégories à valeurs dans une
catégorie telle (Cat), on trouve que la catégorie \inserted[typed]{(Bicat)} des (Cat)-catégories
est équivalente à la sous-catégorie pleine de la catégorie des ensembles
bi-semisimpliciaux

\begin{diagramblock}\[\begin{array}{cccc}
X^{\deleted{00}}_{\hand{00}} & X^{\deleted{01}}_{\hand{01}} & \ldots & X^{\deleted{0n}}_{\hand{0n}}\\
X^{\deleted{10}}_{\hand{10}} & X^{\deleted{11}}_{\hand{11}} & \ldots & X^{\deleted{1n}}_{\hand{1n}}\\
\ldots\ldots\ldots & & &
\end{array}\]\end{diagramblock}

formée des objets qui ligne par ligne et colonne par colonne satisfoht la
condition d'exctitude familière (transformer sommes amalgamées en produit
fibrés). De même la catégorie \xout{des} (Tricat) des (Bicat)-catégories
\unsure{s's'interprète|s'interprète} en t ermes d'objets tri-\inserted[typed]{semi}simpliciaux, etc.

Exemple: soit C une catégorie, posons $C^{\deleted{o}}_{\hand{0}}$=C, $C^{\deleted{1}}_{\hand{1}}$=$\underline{Fl}$(C), avec structure
$p_{1}$,$p_{2}$, évidentes (d'où une (Cat)-catétorie dont les composantes semi-sim
pliciales supérieures sont les catégories $\underline{Fl}^{n}$(C) des n-flêches de C ...)
L'ensemble bi-semisimplicial correspondant a pour composantes les ensembl
Hom($\Delta^{n}\times\Delta^{m}$,C). Par exemple, les objets de $X^{\deleted{11}}_{\hand{11}}$ sont les carrés commutatif
de flêches de C, qu'on peut interpréter(de deux manières différentes) com

\margin{6}

\clearpage
% Tapis de Quillen, with Quillen's offprint and preprint (cote 162-5, typed part): archive cote 162-5, site page 7 (PDF page 8).
% typescript. Draft compared: none.
% AI-transcribed by Claude Opus 5.5 (claude-opus-5-5, adaptive thinking, effort high), Message Batches API, batch t162_full (msgbatch_01UAh3zyWPd454AeDkhnmQ33), 2026-11-13; not yet checked by eye.
\scanpage{162-5/7}
\mspage{\hand{6}}

morphismes de flêches.

Notons une symmétrie évidente dans la c atégorie (Bicat), interprété
en termes d\xout{e}'ensemble bi-semisimpliciaux, obtenue en échangeant les lign
et les colonnes dqns \xout{\illegible{1}} de tels ensembles bi... Un ensemble bisemisimpl
al qui ... peut donc de deux façons différentes être considéré comme déf
nissant une bicatégorie. De même le groupe symétrique $\mathrm{Sym}_n$ opère sur
la catégorie \xout{de} (nupcat) des catégories n-uples ..

Une 2-catégorie peut être définie comme constituée par un ensemble
\xout{discret} $C_0^{\hand{0}}$ (l'ensemble des objets ), qu'on interprètera comme une catégo
rie discrète, plus la donnée pour tout $x,y \hand{$\in$} C_0^{\hand{0}}$ d'une catégorie $\underline{\mathrm{Hom}}(x,y)$,
plus \xout{\illegible{1}} \inserted[typed]{des foncteurs} composition entre celles-ci ... Prenant la catégorie somme \xout{\illegible{1}} C
des $\underline{\mathrm{Hom}}(x,y)$, On peut exprimer la donnée de ces foncteurs de composition
et des sources et buts\inserted{\unsure{$x,y$ )}} dans les catégories $\underline{\mathrm{Hom}}(x,y)$, comme définis par
des foncteurs "source" et "but" $C_{1}^{\hand{1}} \to C_{0}^{\hand{0}}$, et un foncteur de composition
$C_{1}^{\hand{1}} \times_{C_0^{\hand{0}}} C_{1}^{\hand{1}} \longrightarrow C_1^{\hand{1}}$ , les axiomes des 2-catégories correspondant exactement au
axiomes des(Cat)-catégories. De cette façon, la catégorie des 2-catégor
est équivalente à la catégorie des (Cat)-catégories qui ont une "catégo
d'objets" qui est discrète, ou encore à la catégorie des Bicatégories
qui, interprétées comme ensembles bi-semisimpliciaux, ont \xout{une} comme
première ligne un ensemble semisimplicial discret (i.e. un foncteur cons
tant sur \hand{$\Delta$} ). Il parait qu'on a une interprétation simple dans le même
esprit de la notion de n-catégorie en termes d'ensembles n-semi-simplici
\margin{$\|$} \xout{au} Demander à Quillen relations entre la notion de m-Catégorie dans
\margin{$\|$} la catégorie (n-Cat), et celle d\over{e}{'} ensemble (m+n)-semisimplicial.
\margin{$\|$} \xout{(en particulier entre cette dernière notion et celle de}

L'interprétation standard semisimpliciale permet d'interpréter

\margin{7}

\clearpage
% Tapis de Quillen, with Quillen's offprint and preprint (cote 162-5, typed part): archive cote 162-5, site page 8 (PDF page 9).
% typescript. Draft compared: none.
% AI-transcribed by Claude Opus 5.5 (claude-opus-5-5, adaptive thinking, effort high), Message Batches API, batch pilotnext (msgbatch_01A9sdu8SYq23bEtqX9tUncQ), 2026-11-13; not yet checked by eye.
\scanpage{162-5/8}
\margin{\unsure{!} Attention à la catégorie \uline{Hom} \hand{\unsure{7}} grad., qui s'identififie à la catégorie \unsure{définie} \unsure{par un} $R\mathrm{Hom}_{\mathbb{Z}}(C_{\bullet}, C'_{\bullet})$, ce point devrait être élucidé!}

de même les (Cat)-groupes(ou groupes à valeurs dans (Cat)) comme étant
les groupes semisimpliciaux qui "sont"des catégories i.c. qui \xout{sexdéfinid}
\xout{sent} sat\unsure{is}font la condition d'exactitude familière, permettant d'en déte
miner tous les composants à l'aide de $C_{o\xout{m}}$ et $C_1$ . Appliquant le tapis de
Dold, dans le cas commutatif, on trouve que les (Cat)-groupes abéliens
forment une catégorie équivalente à celle des complexes de chaines de
groupes abéliens qui sont de longueur 1:
\[ 0 \leftarrow C_o \leftarrow C_1 \leftarrow 0 \quad . \]
Quand on explicite la (Cat)-groupe abélien associé à un tel complexe de
chaines, on trouve "ma" construction chérie da la "catégorie de Picard"
associée au complexe: objets ceux de $C_o$ , flêches de x dans y les z dans
tels que $dz = y-x$ . On trouve donc à isomorphisme \unsure{près|ppès} tous les group
abéliens de (Cat) ainsi. Un complément important est que toute catégori
\margin{\unsure{Objections déjà} \unsure{mais cela contredit} \unsure{ce qui est dit des} \illegible{2} \unsure{Picard} \illegible{3} \unsure{Mais} \illegible{2} \unsure{de Picard} \illegible{4} \unsure{(Cat)} \illegible{2}}
de Picard \inserted[typed]{(resp. \deleted{et}} commu\deleted{tative)}est équivalente (au sens technique qu'on ne précis
pas ici) à \xout{laxcatégoxxie} un vrai (Cat)-groupe \inserted[typed]{\deleted{(resp. (Cat)-groupe}} \deleted{(commutatif)} . \hand{Quand on}
\hand{\unsure{identifie} mod ``\unsure{Gr}-équivalence'', \unsure{cela} correspond à travailler dans la catégorie \uline{dérivée}}
\hand{\unsure{des complexes}, \unsure{seuls invariants} \unsure{étant} $\pi_0, \pi_1$)}
On notera qu'on a une \xout{lxxxxx} classification beaucoup plus restrein
que pour les H-espaces (resp. les H-espaces homotopiquement comm.) mod
équivalence d'homotopie, de façon précise entre les groupes resp. les gr
oupes abéliens de la catégorie homotopique. La raison en est quon est beau
coup plus restrictif dans le contexte catégorique pour le sens da l'asso
tivité (resp. et de la commutativité), puisqu'on ne demande pas seulemen
l'existence d'isomorphismes d'associ\unsure{ai}tivité\unsure{$_x$} (qui suffirait pour avoir u
groupe dans \inserted[typed]{la catégorie homotopique} (Cat)' $\simeq$ (Ssimpl)', modulo l'existence d'unités et d'invers
mais qu'on suppose \uline{données} ces isomorphismes, de façon à satisfaire à ce
taines compatibilités: c'est ce qui permet, me semble-t-il, de prouver qu
\margin{(}équivalence près, on peut supposer qu'il s'agit de vrais groupes dans
(Cat), \deleted{resp. de vraie commutativité.}

\mspage{8}

\clearpage
% Tapis de Quillen, with Quillen's offprint and preprint (cote 162-5, typed part): archive cote 162-5, site page 9 (PDF page 10).
% typescript. Draft compared: none.
% AI-transcribed by Claude Opus 5.5 (claude-opus-5-5, adaptive thinking, effort high), Message Batches API, batch t162_full (msgbatch_01UAh3zyWPd454AeDkhnmQ33), 2026-11-13; not yet checked by eye.
\scanpage{162-5/9}
\mspage{\hand{8}}

% NOTE: the right edge of the typescript is cut off on the scan; line ends are transcribed as far as visible
On doit trouver une approximation meilleure d'un cran à la classifi
tion topologique en prenant des  (2-Cat)-groupes resp. des (2-Cat)-group
abéliens,  ou plus généralement des " Gr-2-catégories "  resp.
"\unsure{Grab}-2-catégories" . (NB  \unsure{Une} Grab-catégorie défi\xout{\unsure{k}}nit une Grab-2catégor
comme une catégorie définit une 2-catégorie; du point de vue semi-simpli
al, cette opération serait analogue à celle de "cosquélette$^{\text{x}}$ tronqué";
Mais également une Grab-2-catégorie défini\over{e}{t} \over{en}{une} G\unsure{r}ab-catégorie$^{\text{x}}$; savoir la
catégorie des automorphismes de l'objet unité; cette \xout{\illegible{1}} opération serait
analogue à celle de "tuage des groupes d'homotopie en dim \hand{$>$}\unsure{\over{2}{1}}"). Pour l
notion gro\xout{i}ssière ( Gra\xout{b} ou Grab) l'idée de ma lettre à Verdier était qu
doit trouver la même chose que les tronqués \inserted{au cran 2}, dans la catégorie \uline{homotopic}
des H-espaces  -  i.e. ce qu'on déduit des H-espaces par tuage des group
d'homotopie en dimensions \hand{$>$}2 .

\margin{\hand{$\|$} (double vertical line in the left margin, from here to ``lacets'')}(Demander à Quillen relations exactes entre \xout{H}objets-groupes et objet
\inserted[typed]{\xout{\unsure{pointésx}}}connexes de la catégorie homotopique\inserted{ponctuée}: forment-ils des catégories équival
tes par le foncteur espace des lacets\hand{?}).
\margin{\unsure{non, on a} \unsure{un foncteur} \unsure{pl. fidèle} (?) \illegible{3} \unsure{négatifs} \illegible{1}}

Pour les \uline{vrais} groupes \inserted{abéliens}, de (2-\unsure{\over{C}{C}}at), l'approche de Quillen consiste à
les identifier \unsure{à'abord|d'abord} à des groupes abéliens de la catégorie (Bicat),
donc à des grpupes abéliens bi-semiesimpliciaux particuliers (conditions
sur chaque ligne et chaque colonne). Appliquant la version bi-semisimpli
ale du tapis de Dold, on trouve que \xout{\illegible{1}} la catégorie de ces objets
est\xout{h} équivalente à la catégorie des \xout{objets} bicomplexes \inserted{de chaines} de groupes  abéli
qui sont nuls en dehors du carré  \xout{D}$0 \leq i,j \leq 1$.  \xout{Quillen dit que} La \deleted{colonn}\inserted{cond.}
que la première ligne corresponde à un \xout{ensemble} groupe semisimplicial
trivial signifie alors qu'elle n'a que le terme de degré zéro, i.e. que
le bicomplexe est réduit à

\begin{diagramblock}
\begin{tikzcd}
C_{oo} & 0 \arrow[l] & \hand{$0$} \arrow[l, "\hand{}"'] \\
\hand{C}C_{1o} \arrow[u] & C_{11} \arrow[l] \arrow[u, "\hand{}"] & \hand{$0$} \arrow[l, hook', "\hand{}"']
\end{tikzcd}
\end{diagramblock}

\margin{\hand{9}}

\clearpage
% Tapis de Quillen, with Quillen's offprint and preprint (cote 162-5, typed part): archive cote 162-5, site page 10 (PDF page 11).
% typescript. Draft compared: none.
% AI-transcribed by Claude Opus 5.5 (claude-opus-5-5, adaptive thinking, effort high), Message Batches API, batch t162_full (msgbatch_01UAh3zyWPd454AeDkhnmQ33), 2026-11-13; not yet checked by eye.
\scanpage{162-5/10}
\mspage{9}

\unsure{Ces|Cas}\xout{x} bicomplexes s'identifient aux complexes de chhines  de \unsure{longueur} 2.
Donc la catégorie de (2-Cat)-groupes abéliens est équivelente à celle de
complexes de chaines de longueur 2. L'explicitation de la 2-catégorie
associée à un tel complexe de chaines est mon interprétation chérie. Il
faudrait encore expliciter \xout{lesxmorphismes}  2-Grab-catégorie  des
homomorphismes (en un sens à préciser) entre deux\xout{x} tel\xout{l}es  2-Grab-catég
% NOTE: the right edge of the page is cut; the line ends are not visible
ries strictes, en termes d'un complexe RHom.

L'idée qui se dégage de cette approche, et que Quillen on espère v
expl\unsure{i}citer, c'est que la \xout{notionxdexcatégoriexde} notion\xout{s} de groupe \uline{abélie}
dans (n-Cat) correspond à celle de complexe de \inserted{chaines de} longueur n, en raisonnant
soit mod. isomorphisme de part et d'autre, soit mod. une notion naturell
de n-équivalence d'un côté, et dans la catégorie dérivée de l'autre. Cec
doit pouvoir se préciser en utilisant des Hom$^{\bullet}$ resp. des RHom$^{\bullet}$.

Si maintenant C est une catégorie avec un ``produit tensoriel'' assoc
tif (et éventuellement unitaire), i.e. un bifoncteur avec des isomorphis
mes d'associativité (mais sans inverse ...), on peut lui associer \unsure{sans}
doute une  Gr-catégorie universelle (tout comme à un monoïde on associe
groupe symmétrisé), d'où des invariants, qu'on pourra désigner par K$_{o}$(C)
et K$_{1}$(C). On suppose qu'on pourra de même trouver, pour tout n, une
Gr-n-catégorie universelle où envoyer C de façon compatib\xout{b}le avec $\otimes$ (\unsure{tout}
cela dans un sens à préciser), dont les invariants \hand{$\pi_n$}  seraient par déf
tion les K$_{n}$(C). \unsure{Un} cas intéressant est celui où C est une catégorie ave
produits finis (par exemple une catégorie additive) et $\otimes$ est le produit\unsure{;}
\xout{mais}\over{(}{[}Mais on notera que dans le cas d'une catégorie non seulement additi
mais abélienne, il y a une notion plus adéquate du K$_{o}$ , qui devrait corr
pondre à une notion plus adéquate de la suite des K$_{i}$ , inspirée du tapis
de Roos; sans doute c'est en termes de (``vraie'') catégories triangulées

\margin{10}

\clearpage
% Tapis de Quillen, with Quillen's offprint and preprint (cote 162-5, typed part): archive cote 162-5, site page 11 (PDF page 12).
% typescript. Draft compared: none.
% AI-transcribed by Claude Opus 5.5 (claude-opus-5-5, adaptive thinking, effort high), Message Batches API, batch t162_full (msgbatch_01UAh3zyWPd454AeDkhnmQ33), 2026-11-13; not yet checked by eye.
\scanpage{162-5/11}
\mspage{10}

qu'i\unsure{l} conviendra de poser les définitions adéquates). - Peut-être Quill
\margin{\unsure{sic}}a-t-il une définition semi-simpliciale plus directe, sans passer par des
n-catégories, des $K_i$ et plus précisément du H-espace qui devrait correspo
\inserted[typed]{en général}
dre à C. \xout{\illegible{2}} Ce n'est certainement pas l'ensemble semi-simplicial
SC \xout{\unsure{engendrés}} associé à C, puisque dans le cas d'une catégorie additive,
C est homotopiquement trivial (en effet le foncteur identique est ``homoto
au foncteur nul, puisqu'il suffit d'avoir un homomorphisme d'un foncteur
dans l'autre; d'où une équivalence d'homotopie entre C et la catégorie
nulle). Pourtant, il semble bien que dans le cas où C est une Gr-catégor
\margin{\unsure{?}}i.e. admet des inverses, alors ce soit bien la construction ``naïve'' SC qu
est la bonne, puisqu'on trouve bien un $\hand{\pi}_0$ et un $\hand{\pi}_1$ qui sont aussi les
$\hand{\pi}_0$ et $\hand{\pi}_1$ au sens catégorie sans Gr-structure.

\hand{\unsure{À expliquer pour} des Gr-n-catégories, les $\pi_i$ au sens}
\margin{\unsure{?}}\hand{semi-simplicial sont aussi les $\pi_i$ au sens Gr-n-catégories,}
\hand{savoir les objets à isom près dans divers \uline{Hom}$(1,1)$\,..... Il y a}
\hand{aussi une interprétation \unsure{de l'opération} \deleted{\unsure{Steenrod} \illegible{2}}}
\hand{\deleted{\unsure{constructions d'opérations des shifts géométriques évidentes}} sur les}
\hand{n-catégories.....}

\margin{11}

\clearpage
% Tapis de Quillen, with Quillen's offprint and preprint (cote 162-5, typed part): archive cote 162-5, site page 12 (PDF page 13).
% typescript. Draft compared: none.
% AI-transcribed by Claude Opus 5.5 (claude-opus-5-5, adaptive thinking, effort high), Message Batches API, batch t162_full (msgbatch_01UAh3zyWPd454AeDkhnmQ33), 2026-11-13; not yet checked by eye.
\scanpage{162-5/12}
\mspage{11}

III Point de vue ``motivique'' en théorie du cobordisme.

% NOTE: the right edge of the scan cuts off the ends of several typed lines; the visible text is transcribed as it stands
Soit C la catégorie des variétés différebtiables\inserted{(\unsure{compactes?})} (pas nécessairem
orinhtables), \inserted[typed]{\uline{les morphismes étant les classes d'homotopie d'applications continues}} Si B est une catégorie, on s'intéresse aux couples

$(F.,F^{\cdot})$ d'un foncteur covariant et d'un foncteur contravariant de C

dans B, satisfaisant les conditions que pour tout $X\in \mathrm{Ob}\, C$, on a $F.(X)=$

$F^{\cdot}(X)$, et que si on a un produit fibré ordinaire
\begin{diagramblock}\begin{tikzcd}
 & X' \arrow[dl,"g'"'] \arrow[dr,"f'"] & \\
X \arrow[dr,"f"'] & & Y' \arrow[dl,"g"] \\
 & Y &
\end{tikzcd}\end{diagramblock}
,

avec f et g transversaux, alors on a $g^{\cdot}f. = f!\, g'^{\cdot}$\inserted[typed]{(on pose $f.=F.(f)$,$f^{\cdot}=F^{\cdot}($}. La question est

de trouver une B et un couple $(F.,F^{\cdot})$ qui soient ``universels'' dans u

sens évident. Quillen prouve que la construction suivante donne une

telle solution.

Pour toute variété $X\in \mathrm{ob}\, C$ , on définit un anneau commutatif

B(X) (l'anneau de bordisme de X) en prenant comme éléments les class

\xout{de} d'éléments de C au dessus de X, pour la relation de cobordisme:

$f'{:}X'\to X$ et $f''{:}X''\to X$ sont équivalents, s'il existe une variété à bord

Z, et un $f{:}Z\to X$, tel que bord $Z = X'\amalg X''$ et $f'$ et $f''$ sont induits

par f. La somme de B(X) est la somme disjointe (c'est une loi de grou

\xout{\unsure{inverse étant}} où tout élément est d'ordre \hand{\unsure{2}} \xout{\illegible{1}}2), le produit est

donné par le produit fibré, en utilisant le théorème de trahsversalit

de Thom pour bouger un peu f'', de sorte à le rendre transversal à f'.
\margin{NB L'anneau est de plus \unsure{gradué} par la dimension relative ....}
Si on a $f{:}\ X\to Y$, on définit $f. : B(X)\to B(Y)$ de façon évidente (tou

variété sur X est sur Y), et $f^{\cdot}{:}\ B(Y)\to B(X)$ par produit fibré (uti

sant encore le th. de transversalité). De cette façon, B(X) devient u

foncteur-anneau \inserted{(\unsure{gradué})} contravariant en X, et un foncteur groupe covariant e

\mspage{12}

\clearpage
% Tapis de Quillen, with Quillen's offprint and preprint (cote 162-5, typed part): archive cote 162-5, site page 13 (PDF page 14).
% typescript. Draft compared: none.
% AI-transcribed by Claude Opus 5.5 (claude-opus-5-5, adaptive thinking, effort high), Message Batches API, batch t162_full (msgbatch_01UAh3zyWPd454AeDkhnmQ33), 2026-11-13; not yet checked by eye.
\scanpage{162-5/13}
\mspage{12}

X. De plus, on a la formule\inserted{\unsure{du type $g'f_{.} = f'_{.}g'$} ci-dessus, impliquant la formule} de projection habituelle. Ceci posé, les
constructions bien connues en théorie des correspondances permettent
de définir une catégorie $\underline{B}$ dont lds objets sont ceux de C, et les
Hom étant donnés par
\[ \mathrm{Hom}_{\underline{B}}(X,Y) = B(XxY) \quad . \]
On a un isomorphisme canonique $\underline{B} \simeq \underline{B}^{o}$ , qui est l'identité sur les
objets et la symmétrie $B(XxY)\simeq B(YxX)$ sur les flêches, et un foncteur
évident \xout{$\underline{E}$xxx$\underline{E}$x} $F_{.}$ : $C \to \underline{B}$ , qui est l'identité sur les objets et don
% NOTE: the right edge of the page is cut; the line ends with "don" (continuation not visible)
par les graphes sur les flêches. Composé avec la symmétrie précédente
on trouve un foncteur $F^{\cdot}$: \xout{Ex} $C \to \underline{B}$. On vérifie enfin que le couple
($F_{.}$,$F^{\cdot}$) satisfait aux conditions envisagées dans l'alinéa 1 (ceci aus
% NOTE: right edge cut; "aus" (continuation not visible)
devrait être formalisé au cas des correspondances génér ales ...). \unsure{On}
vérifie enfin qu'il a la propriété universelle qu'il faut. Lorsqu'on
se donne un couple ($F'_{!}$,$F'^{\cdot}$) à val eursdans une catégorie B', il faut
trouver une factorisation par B. Elle est toute trouvée sur les
objets, il faut la définir sur les flêches, donc définir, pour X,Y
dans C, une application
\[ (*) \qquad B(XxY) \to \mathrm{Hom}(X',Y') \quad . \]
Mais si Z est auddessus de XxY, d'où $p:Z \to X$ et $q:Z \to Y$, on lui assoc
% NOTE: right edge cut; "assoc" (continuation not visible)
un élément du deuxième membre de (*) comme le composé
\[ X' \xrightarrow{F'^{\cdot}(p)} Z' \xrightarrow{F'_{!}(\unsure{q})} Y' \quad . \]
Il faut prouver que l'élément obtenu en \unsure{remplaçant|remphçant} Z\hand{\unsure{,}} par un $Z_1$ cobord
% NOTE: right edge cut; "cobord" (continuation not visible); small handwritten marks ". 2. ," above "Z par"
dant au dessus de X\inserted{\unsure{ne varie pas|n'a pas varié}}. Or si $g$: $T \to XxY$ réalise le cobordisme, on
voit (regardant seulement T et la décomposition bord(T) = $Z \amalg Z_1$)
qu'il existe une variété $\overline{T}$ sur $S^1$ (obtenue en ``doublant'' T) \xout{etxxxx}
\xout{applicationxTxxxxx} trqnsversale au dessus d'un couple de poihts $s_o$,s
% NOTE: right edge cut; "$s_o$,s" (index of the second point not visible)
de $S^1$, tels que $Z_i$ soit l'image inverse \inserted{dans $\overline{T}$} de $s_i$ (i=0,1); d'autre part,

\margin{13}

\clearpage
% Tapis de Quillen, with Quillen's offprint and preprint (cote 162-5, typed part): archive cote 162-5, site page 14 (PDF page 15).
% typescript. Draft compared: none.
% AI-transcribed by Claude Opus 5.5 (claude-opus-5-5, adaptive thinking, effort high), Message Batches API, batch t162_full (msgbatch_01UAh3zyWPd454AeDkhnmQ33), 2026-11-13; not yet checked by eye.
\scanpage{162-5/14}
\mspage{13}

g manifestement se prolonge en $\overline{g}$ : $\overline{T}\to X\times Y$ . Cela nous ramène à prouv
\inserted{le composé}
que \hand{(}$\unsure{\overline{T}'}\xrightarrow{\hand{\alpha}}\unsure{\overline{Z}'_i}\xrightarrow{\hand{\unsure{\alpha_i}}}\overline{T}'$\hand{)} ne dépend pas de i=0,1. Ecrivons X au lieu de $\overline{T}'$,

considérons le diagramme cartésien d'inclusions
\begin{diagramblock}\begin{tikzcd}
X_s \arrow[r, "\hand{\alpha_s}"] \arrow[d, "\hand{\alpha_s}"'] & X \arrow[d, "\hand{\unsure{\Gamma_{\underline{pr}_1(x)}}}"] \\
X \arrow[r, "i_s"] & S^1\times X
\end{tikzcd}\end{diagramblock}
\[ , \]

on a alors $\hand{\alpha}^{\hand{!}}_s\,\hand{\alpha}_{s*} = \hand{\Gamma}^{\hand{!}}i_{s\hand{\unsure{*}}}$ , qui ne dépend pas de s puisque $i_s$ n'en

dépend pas à homotopie près. - Le reste est sans doute l'AQT.

Sur la strcture des B.(X). Ce sont tous des algèbres graduées

sur B.(pt)=\hand{$\Lambda$} , qui est une algèbre graduée à degrés positifs, réduit
\margin{\illegible{1}}
à $\underline{F}_2$ en degré\xout{z} 0. C'est d'ailleurs une algèbre de polynômes connue\inserted{(à une infinité de générateurs \unsure{$X_i$} \illegible{2})}
\inserted[typed]{d'anneauxpgradués}
On a un homomorphisme évident B.(X) $\to H^{\hand{\unsure{B}}}(X)$, nul en degrés \hand{$>$}0, qui
\margin{NB $H^*(X)$ $(=H^*(X,\underline{F}_2))$}
à tout Z sur X associe l'image de 1 \hand{$\in$} $H^0(Z)$ par l'homomorphisme de Gys
\margin{\unsure{en fait} $H(X)\simeq B.(X)/\Lambda^+B.(X)$}
associé à Z$\to$ X. Cet homomorphisme par Thom est surjectif (toute cla

se de cohomologie provient d'une image de variété compacte....); \xout{xxte}

On ppouve même qu'il existe un relèvement fonctoriel en X
\margin{?}
$H^\bullet$(X) $\to$ B.(X), \hand{(}compatible avec multiplication\hand{)}\inserted{il est d'ailleurs \unsure{unique s'il existe}} \xout{xxx}tel que l'on ait:

\unsure{a)} L'homomorphisme correspondant $H^\bullet(X)\otimes_{\underline{F}_2}\Lambda\simeq$ B.(X) est un iso
\margin{?}
morphisme (ce sera vrai pour tout relèvement individuellement ...), e

b) \hand{(}\deleted{Le\xout{xx}}\inserted{\unsure{quelconque}} relèvement est compatible aussi avec les hom. de Gysin

Ainsi, la catégorie $\underline{B}$ admet un foncteur pleinement fidèle dans la

catégorie des modules gradués libres de type fini sur \hand{$\Lambda$}, les modules

de type fini qu'on obtient ainsi ayant des générateurs à degrés \hand{$<$} 0

satisfaisant (pour X connexe i.e. quand il y a un seul générateur de
\margin{i.e. l'\unsure{enveloppe} \unsure{``Karoubienne''}}
degré zéro\hand{)} la dualité de Poincaré. Quand on prend la catégorie \uline{quas
\uline{abélienne}\inserted{\unsure{B}}-associée à $\underline{B}$, on trouve donc tous les modules gradués libre

de type fini à générateurs \xout{négatifsx} en degrés \hand{$\leq$} 0 seulement. C'est u

\margin{14}

% COMPILE FIX (combined file only): 1 closing brace(s) added; the page's LaTeX leaves them open
}

\clearpage
% Tapis de Quillen, with Quillen's offprint and preprint (cote 162-5, typed part): archive cote 162-5, site page 15 (PDF page 16).
% typescript. Draft compared: none.
% AI-transcribed by Claude Opus 5.5 (claude-opus-5-5, adaptive thinking, effort high), Message Batches API, batch t162_full (msgbatch_01UAh3zyWPd454AeDkhnmQ33), 2026-11-13; not yet checked by eye.
\scanpage{162-5/15}
\mspage{14}

\inserted[typed]{additive(et même \margin{=karoubienne+additive} quasi-abélienne)}
catégorie \inserted[typed]{additive(et même quasi-abélienne)} avec $\otimes$ qui a toutes les vertus, sauf d'être abélienne, et
sauf d"admettre des \uline{Hom} internes. Pour trouver ces derniers, on peu
si on veut paraphraser la construction de la catégorie des motifs,
\xout{\unsure{introduisant un motif}} \inserted{en} considérant l'objet \inserted{$\hand{\Lambda}(-1)$} de $\overset{\hand{$\sim$}}{\underline{B}}$ conoyau de
$p^{\hand{$*$}} : e \longrightarrow S$, où $p:S^1 \longrightarrow e$ est l'application du cercle sur un
point. C'est un objet de poids 1 (le poids est défini par le degré
des générateurs, changés de signe). On constate aussitôt que le
foncteur $M \mapsto M(-1) = M\otimes \hand{$\Lambda$}(-1)$ de $\overset{\hand{$\sim$}}{\underline{B}}$ dans lui-même est pleinement fidèl
ce qui permet alors, par une construction formelle standard, de com
\margin{\unsure{$\ell\varphi$}} ter $\overset{\hand{$\sim$}}{\underline{B}}$ en une catégorie où $\hand{$\Lambda$}(-1)$ devient inversible. Il se trouve
alors que dans cette dernière les \uline{Hom}(X,Y) internes existent et
\xout{existent} sont isomorphes à $\overset{\hand{$\vee$}}{X}\otimes Y$, ... en fait , on trouve maintenant\hand{\unsure{(|l}}
catégorie de\xout{s} \uline{tous} les modules gradués libres de type fini sur $\hand{$\Lambda$}$.
Bien sûr, cette catégorie n'est pas encore abélienne pour autant,
puisque $\hand{$\Lambda$}$ n'est pas un corps; \xout{pour} \inserted[typed]{si on veut} définir la catégorie abélienne
engendrée, on s'attend à trouver la\xout{s} catégorie desmodules de \deleted{type} \inserted{\unsure{présentation}} f
\margin{\unsure{$\S$}} sur $\hand{$\Lambda$}$.( \inserted{\unsure{qui est cohérent bien que $\Lambda$ ne soit pas noethérien, car $\Lambda$ est}} \deleted{est-il noethérien, i.e. le nombre de générateurs est-il
fini ? )}\hand{, \unsure{lim} d'anneaux \unsure{noeth.} à \unsure{hom.} de transition \uline{plats} )}.

Cette construction ressemblerait à celle d'une théorie des
motifs où on n'aurait pas passé à l'équivalence homologique pour
les cycles, mais où on aurait gardé l'équivalence rationelle, - qui
est bien trop fine \xout{\unsure{par}}à beaucoup d'égards. Pour trouver un équiva
lent plus étroit, il convinnt de reprendre la théorie prédédente,on
pourrait formuler le même Pb universel que plus haut, mais en prena
l'équivalence homologique (via le graphe) des morphismes de variété
au lieu de l'équivalence homotopique. \deleted{Mais alors il semble que l'on
trouve comme catégorie solution celle dont les objets sont les vari}

\mspage{15}

\clearpage
% Tapis de Quillen, with Quillen's offprint and preprint (cote 162-5, typed part): archive cote 162-5, site page 16 (PDF page 17).
% typescript. Draft compared: none.
% AI-transcribed by Claude Opus 5.5 (claude-opus-5-5, adaptive thinking, effort high), Message Batches API, batch t162_full (msgbatch_01UAh3zyWPd454AeDkhnmQ33), 2026-11-13; not yet checked by eye.
\scanpage{162-5/16}
\mspage{15}

% NOTE: the right edge of the typescript is cut off on the scan; line ends are transcribed as far as visible
\deleted{tés, les morphismes \uline{toutes} les co\unsure{r}respondances \inserted[typed]{co}homologiques (en ve
du thorème de Thom rappelé plus haut). On trouve donc un foncteur
pleinement fidèle de cette catégorie dans celle des vectoriels
gradués \xout{dxxd} positivement de dim. finie sur $\underline{F}_2$ , qui est déjà abéli
ne, et si on veut rendre comme plus haut $\underline{F}_2(-1)$ inversible, on trou
tous les vectoriels gradués de dim. finie sur $\underline{F}_2$. C'est un peu tro
trivial comme structure, et on peut faire mieux,} \xout{exprenantxcomme}
\inserted[typed]{Pour trouver une catégorie universelle,}
\xout{correspondancesxPourxceci} on associe à toute X \xout{\illegible{2}}
\xout{lx} l'ensemble \inserted[typed]{S(X)} des sous-variétés \xout{S(X)} de X, modulo équivalence homo
gique. Donc S(X) $\subset$ H$^{\cdot}$(X), et on prouve qu'on trouve en fait un sous
anneau de H$^{\cdot}$(X) (grâce à un théorème de Thom qui caractérise S(X) \hand{en}
termes d'opérations de Steenrod, cf plus bas). On prouve de même
(par le même théorème) que les S(XxY) sont stables par composition
correspondances cohomologiques( cela ne semble pas trivial, puisque
les S(X) ne sont pas stables par Gysin...). D'où une catégorie add
tive dont les objets sont ceux de C, les flêches sont les \inserted{él. des} S(XxY),
Il semble qu'elle doive jouer le rôle universel que j'ai envisagé
plus haut. En tous cas, sa définition géométrique est fort simple
Nous allons en donner une interpétation algébrique. Pour ceci,
considérons l'anneau des opérations cohomologiques \inserted[typed]{``stables''} \deleted{additives sur l}
\deleted{des endomorphismes du}
\margin{\unsure{?}}\xout{lx} \deleted{foncteur X --- H$^{\cdot}$(X) qui sont additifs, et qui en degré zéro se
réduisent à un multiple de l'identité}\hand{/} C'est un anneau gradué, avec
application diagonale anticommutative, \inserted[typed]{i.e. commutative (car. 2 !)} associative et unitaire. Son
dual gradué est une algèbre \inserted[typed]{graduée anti\xout{o}commutative donc \xout{commutative}} commutative,avec application diagonale
associative et unitaire; il est réduit à $\underline{F}_2$ en degré zéro. Donc il
corespond à un schéma en groupes \unsure{affine} G sur $\underline{F}_2$ , qui est d'aille
le foncteur des \unsure{$\otimes$}-automorphismes \inserted[typed]{additifs} du foncteur H$^{\cdot}$ sur C \hand{qui opèrent
trivialement sur \deleted{\illegible{1}} \unsure{$H^2(-1)$|$\underline{F}_2(-1)$}. \deleted{\illegible{1}}}

\margin{16}

\clearpage
% Tapis de Quillen, with Quillen's offprint and preprint (cote 162-5, typed part): archive cote 162-5, site page 17 (PDF page 18).
% typescript. Draft compared: none.
% AI-transcribed by Claude Opus 5.5 (claude-opus-5-5, adaptive thinking, effort high), Message Batches API, batch t162_full (msgbatch_01UAh3zyWPd454AeDkhnmQ33), 2026-11-13; not yet checked by eye.
\scanpage{162-5/17}
\mspage{16}

Le théorème de Thom est alors qu'un $x \in H^{\cdot}(X)$ est dans $S(X)$ sss il est
\xout{\illegible{1}}\inserted[typed]{fixe} par G, i.e. satisfait $A^{+}.\ x = 0$. (Pourquoi cette condition est-el%
% NOTE: the right edge of the scan cuts the ends of several typed lines
le déjà nécessaire ?) Une autre façon de le voir est de dire que le
foncteur $H^{\cdot}$ considéré comme allant de la catégorie $\underline{S}$ précédente dans
la catégorie des \xout{\unsure{vectoriels}} représentations de G dans des vectoriels
sur $\underline{F}_2$ , est pleinement fidèle. \xout{Comme G est} Quand on prend l'enveloppe
pseudo-abélienne, on trouve une sous-catégorie pleine de la catégorie des
représentations de G, mais il est douteux qu'on trouve une catégorie abé-%
lienne, voire \unsure{la|ta} catégorie de toutes les représentations de G de dimensi
finie sur $\underline{F}_2$ . (Quelle est la structure de G, qui est prolisse; d'après
\margin{\illegible{20}}
Quillen, il serait pro-\unsure{nilpotent} ....). Ici G bien sûr joue le rôle d'u
groupe de Galois motivique, mais contrairement à la situation familière,
il semble qu'il soit fortement non réductif (et d'ailleurs le corps de
\hand{$\triangleright$} \unsure{définition} dans tout ceci est $\underline{F}_2$ , et non $\underline{Q}$ ! ).

Quillen\deleted{\unsure{t}} donne une généralisation de la construction d'une caté-%
gorie universelle $\underline{B}$ au cas où on travaille avec des variétés pas nécessa
rement propres (mais $f_{.}$ n'étant supposé défini que pour f propre), et
quand on introduit une condition d'orientation \inserted[typed]{stable} sur les \uline{morphismes}
envisagés , via un morphisme $H \to BO(\infty)$, où H est un H-espace (commuta
tif, associatif et unitaire ?): une orientation de $f:X \to Y$ sera un
relèvement homotopique de $X \to BO$, défini par \xout{\illegible{1}} $T_X - f^{\cdot}(T_Y)$ , en $X \to H$
les éléments de $H(X,Y)$ seront définis par des diagrammes de morphismes
orientés $Z \to X$, $Z \to Y$, le deuxième étant propre, mudulo bords de tels.
En tant qu'ensemble, il se définit aussi en termes de la vari\hand{é}té XxY
munie simplement de la famille des supports formée des ensembles fermés
propres sur Y ..... \hand{Les $B^{H}_{m}(X)$ (X \unsure{compact disons}) peuvent}
\hand{se définir homotopiquement à la Thom, comme les groupes d'homotopie}
\hand{\unsure{stables d'espaces} \uline{Hom}$(X, MH(N-m))$ -----.}

\hand{cf exposé. Karoubi à \uline{Bourbaki}}

\margin{17}

\clearpage
% Tapis de Quillen, with Quillen's offprint and preprint (cote 162-5, typed part): archive cote 162-5, site page 32 (PDF page 33).
% typescript. Draft compared: none.
% AI-transcribed by Claude Opus 5.5 (claude-opus-5-5, adaptive thinking, effort high), Message Batches API, batch t162_full (msgbatch_01UAh3zyWPd454AeDkhnmQ33), 2026-11-13; not yet checked by eye.
\scanpage{162-5/32}
\margin{notes du 10.9.68}

Tapis de Quillen.

I \quad \uline{Relations entre catégories et ensembles \unsure{semi-simplicaux}.}

A toute catégorie C,  on associe un ensemble semi-simplicial S(C),
trouvant ainsi un foncteur pleinement fidèle \xout{dexCat}
\[ S; \ (\mathrm{Cat}) \longrightarrow (\mathrm{Ssimpl}) \quad . \]
\xout{Si C est associé à} Les systêmes locaux d'ensemble sur SC correspondent
aux \unsure{fancteurs|foncteurs} sur C qui transforment toute flêche en isomorphisme (i.e.
qui se factorisent par le groupoïde associé à  C). Les $H^i$ sur SC d'un tel
système local ($H^0$ pour ens, $H^1$ pour groupes, $H^i$ quelconques pour groupes
abéliens) s'interprètent en termes des foncteurs $\underleftarrow{\mathrm{lim}}^{(i)}$ \unsure{dérivés|dérigés} de $\underleftarrow{\mathrm{lim}}$ ,
ou si on préfère, des $H^i($ du \uline{topos} C . On voit ainsi à quelle condition
un foncteur  $C \to C'$ induit un\xout{\unsure{e}}  homotopisme  $SC \to SC'$ :  \xout{\unsure{e}}en vertu du
critère cohomologique de Artin-Mazur, il f et s que  pour tout système
de coéfficients  F'sur C', l'homomorphisme naturel $\underset{C'}{\underleftarrow{\mathrm{lim}}}^{(i)}F' \longrightarrow \underset{C}{\underleftarrow{\mathrm{lim}}}^{(i\ )}F$
soit un isomorphisme (pour les i pour lesquels cela a un sens). \xout{\unsure{En}\illegible{1}}

A C on peut associer le topos $\overset{\hand{\sim}}{C}$ , qui varie de façon \uline{covariante} avec C.
(NB le foncteur $C \mapsto \overset{\hand{\sim}}{C}$  n'a plus rien de pleinement fidèle, semble-til ??)
Les systêmes de coéfficients \inserted[typed]{ensemblistes} sur C ( = \inserted[typed]{\unsure{déf}} les foncteurs \inserted[typed]{$C^o \longrightarrow \mathrm{Ens}$} transformant isomor-%
phismes en isomorphismes) correspondent aux faisceaux localement constants,
i.e. les objets localement constants de $\overset{\hand{\sim}}{C}$, définis intrinsèquement en termes
de $\overset{\hand{\sim}}{C}$ .  Ainsi, le fait pour un foncteur $F:C \to C'$ d'induire un\xout{e équivalence}
homotopisme $S(C) \to S(C')$  ne  dépend que du morphisme de topos $\overset{\hand{\sim}}{F}: \overset{\hand{\sim}}{C} \to \overset{\hand{\sim}}{C'}$
induit, et signifie que pour tout faisceau localement constant F' sur C' i.
e. sur $\overset{\hand{\sim}}{C'}$ ,  les applications induites  $H^i(\overset{\hand{\sim}}{C'},F') \longrightarrow H^i(\overset{\hand{\sim}}{C},\overset{\hand{\sim}}{F}{}^{\hand{*}}(F'))$
sont des isomorphismes (pour les i pour lesquels cela a un\xout{xx} sens).

\mspage{32}

\clearpage
% Tapis de Quillen, with Quillen's offprint and preprint (cote 162-5, typed part): archive cote 162-5, site page 33 (PDF page 34).
% typescript. Draft compared: none.
% AI-transcribed by Claude Opus 5.5 (claude-opus-5-5, adaptive thinking, effort high), Message Batches API, batch t162_full (msgbatch_01UAh3zyWPd454AeDkhnmQ33), 2026-11-13; not yet checked by eye.
\scanpage{162-5/33}
On a aussi un foncteur évident
\[ T: (\mathrm{Ssimpl}) \longrightarrow (\mathrm{Cat}) , \]
en associant à tout ensemble semi-simplicial X la catégorie $T(X) = \Delta_{/X}$
des simplexes sur X, dont l'ensemble des objets est la réunion disjointe
des $X_n$ .... (c'est une catégorie fibrée sur la catégorie $\Delta$ des simplexes
types, à fibres les catégories discrètes définies par les $X_n$). Ceci posé,
Quillen prouve que pour tout X, ST(X) est isomorphe canoniquement à X dans
la catégorie homotopique construite avec (Ssimpl), et que pour toute C,
la catégorie TS(C) est \inserted[typed]{canoniquement} ``homotopiquement équivalente à C'' \xout{dans la \illegible{1}} i.e.
\inserted[typed]{canoniquement} isomorphe à C dans la catégorie quotient de (Cat) obtenue en inversant les
foncteurs qui sont des homotopismes. \xout{Il en résulte} Ces isomorphismes sont
fonctoriels en X. Il en résulte formellement qu'un morphisme $f:X\to Y$ dans
(Ssimpl) est un homotopisme si et seulement si il en est ainsi de $T(f):T(X)$
$\to T(Y)$, d'où des foncteurs $S'$: $(\mathrm{Cat})' \longrightarrow (\mathrm{Ssimpl})'$ et $T'$:$(\mathrm{Ssimpl})'\longrightarrow(\mathrm{Cat})'$
entre les catégories ``homotopiques'', \inserted[typed]{construites avec (Cat) resp.(Ssimpl),} \hand{/} qui sont quasi-inverses l'un de l'aut-%
re.

De plus, Quillen construit un isomorphisme canonique et fonctoriel
dans (Cat)' entre C et la catégorie opposée $C^{\mathrm{o}}$ , ou ce qui revient au
même, un isomorphisme canonique et fonctoriel dans (Ssimpl)' entre S(C) et
S($C^{\mathrm{o}}$). \xout{Cela lui permet alors, pour} La définition est telle que le
foncteur induit sur les systêmes locaux sur C transforme le foncteur
contravariant F sur C, transformant toute flêche en flêche inversible, en
le foncteur covariant (i.e. contravariant sur $C^{\mathrm{o}}$) ayant mêmes valeurs sur
les objets, et obtenu sur les flêches en remplaçant F(u) par $F(u)^{-1}$; en
d'autres termes, l'effet de l'homotopisme de Quillen sur les groupoides
fondamentaux est l'isomorphisme évident entre les groupoides fondamentaux
de C et de $C^{\mathrm{o}}$ ,\xout{qui est \unsure{l'identité}} compte tenu que le deuxième est l'opposé du

\mspage{\hand{33}}

\clearpage
% Tapis de Quillen, with Quillen's offprint and preprint (cote 162-5, typed part): archive cote 162-5, site page 34 (PDF page 35).
% typescript. Draft compared: none.
% AI-transcribed by Claude Opus 5.5 (claude-opus-5-5, adaptive thinking, effort high), Message Batches API, batch t162_full (msgbatch_01UAh3zyWPd454AeDkhnmQ33), 2026-11-13; not yet checked by eye.
\scanpage{162-5/34}
premier.  Comme applic\xout{t}ation, Quillen obtient une interprétation faisceauti-%
que de la cohomologie d'un ensemble semi-simplicial à coéfficients dans un
système \inserted{local} con\xout{tra}variant \inserted{F} \xout{\unsure{dx}}(défini classiquement par le complexe cosimplicial
\margin{$\prod$ ?} des  $C^n(F) = \coprod_{x\in X_n} F(x)$ ):  on considère le système local contravariant
défini par F, on l'interprète comme un faisceau sur \inserted[typed]{T(X) i.e. objet de} $(\mathrm{Ssimpl})_{/X}$ , et on
prend sa cohomologie. - Cependant, quand F est un système de coéfficien\unsure{ts|ss}
co\unsure{va}riant pas nécessairement local, on n'a tjrs pas d'interp\xout{\unsure{r}}étation de se\unsure{s|n}
groupes de cohomologie  classiques en t ermes faisc\xout{\unsure{e}}autiques; ni, lorsque F
est \unsure{contravariant}, de son homologie, ou inversement de sa cohomologie
faisceautique en termes classiques.

A propos de la notion de foncteur qui est un homotopisme. Quillen
montre qu\xout{e}'un tel foncteur $F: C \hand{\longrightarrow} C'$ induit une équiva\unsure{lence} entre la sous-%
catégorie triangulée \inserted{$D^b_{lc}(C')$} de la catégorie \inserted{dérivée \inserted{bornée} de celle \deleted{des}} des faisceaux abéliens  sur C', dont
les faisceaux de cohomologie sont des systêmes locaux, et la catégorie ana-%
logue pour C; \xout{\unsure{Ilxsxxx}}  et réciproquement. \xout{Pourxlaxpartiexdirecte} On peut
dans cet énoncé introduire aussi n'importe quel anneau de base (à condition
de le supposer $\neq$ 0 dans le cas de la réciproque); \xout{\unsure{constant}} la partie directe
vaut aussi avec un anneau de boefficients pas néc. constant, mais constant
tordu.  Je pense que ce résultat (facile) doit pouvoir se généraliser ainsi:
Soit  f: X \quad X' un morphisme de topos qui  soit tel que pour tout \xout{système}
faisceau loc. const. sur X', f induise un isomorphisme sur les cohomologies
(avec cas non commutatif inclus).  Supposons que \deleted{\xout{\unsure{le topos}} X \unsure{soit}} \inserted{X et $X'$} \uline{localement}
\margin{\unsure{$\circ$}} \margin{$\|$} \uline{homotopiquement trivial}, i.e. que \xout{tout} pour tout entier n\inserted{$\geq 1$}, tout objet U \deleted{de}
\deleted{$X'$} ait un \inserted[typed]{recouvrement} \xout{recouvrement} par des $U_i \hand{\longrightarrow} U$, tels que  a) tout système local sur U
devient constant sur $U_i$ , et toute section sur U devient constante sur $U_i$ ,
et  b)  pour tout \xout{système}  groupe abélien G, les $H^{\hand{j}}(U,G) \longrightarrow H^j(U_i,G)$ sont
nuls pour $1 \hand{\leq} j \hand{\leq} n$.  Alors le foncteur $D^b_{lc}(X') \longrightarrow D^b_{lc}(X)$  induit par f est

\margin{\unsure{$\circ$} Attention, cette condition n'est \unsure{pratiquement} \uline{pas} satisfaite par les schémas avec leur topologie étale (mais bien par \unsure{schémas loc. noeth.} avec top. Zariski).}

\mspage{34}

\clearpage
% Tapis de Quillen, with Quillen's offprint and preprint (cote 162-5, typed part): archive cote 162-5, site page 35 (PDF page 36).
% typescript. Draft compared: none.
% AI-transcribed by Claude Opus 5.5 (claude-opus-5-5, adaptive thinking, effort high), Message Batches API, batch t162_full (msgbatch_01UAh3zyWPd454AeDkhnmQ33), 2026-11-13; not yet checked by eye.
\scanpage{162-5/35}
\xout{pleinementxfidèlex(respx} une équivalence). Même énoncé si on met dans le
coup un \xout{Anneauxinex} systême local d'anneaux sur X'. \unsure{Enfin|Ehfin}, f induit \xout{une}
\inserted[typed]{\xout{foncteurxplxfidxx(respx} une équivalence)}
\xout{isomorphisme} entre la catégorie des coefficients locaux sur C et celle
des coéff. locaux sur C' \xout{(??xpeut-êtrexcelaxdoitxil}\inserted[typed]{\xout{cettexpropriétéxdoit}}\xout{xêtrexmisedansxles}
\xout{hypothèsesxpréliminairesxsurxfx?)x} Principe de démonstration: on commence
par prouver ce dernier résultat, en notant que \xout{dansxlexca} si un topos est
locqalement hom. trivial, il est loca connexe et loc. simplement connexe,
d'où une bonne théorie du $\hand{\pi}_o$ et du $\hand{\pi}_1$ (qui sont ici discrets), et on est
ra\unsure{\over{n}{m}}ené \xout{pour} à un cas particulier du critère d'homotopisme de Artin et
Mazur, savoir un critère \inserted[typed]{cohomologique} pour qu'un homomorhisme de groupes G\hand{$\longrightarrow$}H (ici \xout{l}des
groupes $\hand{\pi}_1$ de X,X') soit un isomorphisme: il doit induire des isomorphis\xout{mx}
mes sur les H$^o$ et H$^1$ (y inclus dans le cas non commutatif...). On prouve
la pleine fidélité en se ramenant par la suite spectrale aux cas \unsure{des|les} Ext$^i$
de coéfficients locaux. Par la suite spectrale encore, cela résultera du
fait suivant: si X\hand{\unsure{$'$}} est \unsure{localement} hom. trivial, alors la catégorie des
faisceaux abéliens loc. constants est stable par \uline{Ext}$^i$ , et le foncteur
M \hand{$\longmapsto$}M$_X$ de\xout{sxfaisceauxxconstants} (Ab) dans X$_{ab}$ commute aux dits Ext$^i$ .
En fait, X et X' étant loc homp\hand{|}triviaux, les conditions suivantes sur f
seront équivalentes:

a) f est\unsure{b} un homotopisme, i.e. induit pour tout systême local (pas
néc commutatif) sur X'un isomorphisme sur les H$^i$\xout{pourxixxx}.

b) \xout{Mêmexchosexave} f induit une équivalence D$^b_{lc}$ (X')\hand{$\longrightarrow$}D$^b_{lc}$(X) .

\unsure{a\hand{$'$})} f induit une équivalence entre la catégorie des systêmes locaux
abéliens sur X' et X, et des isomorphismes sur les H$^i$ correspondants (donc
on ne prend ici que des coéfficients commutatifs).

\margin{?} J'ignore si on peut dans a) se borner aux systêmes locaux commutatifs,
\xout{ixexxdans}
L'équivalence entre a) et b) fournit une première justification ou
\mspage{35} motivation pour définir des types d'homotopie via la catégorie D\hand{$^b_{lc}$(X)},

\clearpage
% Tapis de Quillen, with Quillen's offprint and preprint (cote 162-5, typed part): archive cote 162-5, site page 36 (PDF page 37).
% typescript. Draft compared: none.
% AI-transcribed by Claude Opus 5.5 (claude-opus-5-5, adaptive thinking, effort high), Message Batches API, batch t162_full (msgbatch_01UAh3zyWPd454AeDkhnmQ33), 2026-11-13; not yet checked by eye.
\scanpage{162-5/36}
éventuellement muni de la sous-catégorie pleine de tous les systêmes locau
sur X, et du foncteur cohomologique sur $D^b_{lc}(X)$ à valeurs dans \unsure{ladite}
catégorie, et bien sûr du produit tensoriel (mais alors on sort de $D^b$
pour entrer dans $D^-$ , redactor demerdetur).

II n-catégories, catégories n-uples, et (Gr)-catégories.

Appelons bi-catégorie une catégorie dans(Cat), i.e. un triple
$(C^{\deleted{o}}_{\hand{0}},C^{\deleted{1}}_{\hand{1}},$ \xout{\unsure{xxoù}}$p_1,p_2,\hand{\pi}$ ), où $C^o$ est une catégorie, $p_1,p_2$: $C^{\deleted{1}}_{\hand{1}}\hand{\rightrightarrows}C^{\deleted{o}}_{\hand{0}}$ deux
foncteurs, $\hand{\pi}:(C^{\deleted{1}}_{\hand{1}},p_1)x_{C^o}(C^{\deleted{1}}_{\hand{1}},p_2)\hand{\longrightarrow}C^{\deleted{1}}_{\hand{1}}$ un foncteur, avec les axiomes
habituels. Interprétant (Cat) comme la sous-catégorie pleine de (Ssimpl)
des objets qui .... , et de même pour les catégories à valeurs dans une
catégorie telle (Cat), on trouve que la catégorie \inserted[typed]{(Bicat)} des (Cat)-catégories
est équivalente à la sous-catégorie pleine de la catégorie des ensembles
bi-semisimpliciaux

\begin{diagramblock}
\[
\begin{array}{llll}
X^{\deleted{00}}_{\hand{00}} & X^{\deleted{01}}_{\hand{01}} & \ldots & X^{\deleted{0,n}}_{\hand{0n}}\\
X^{\deleted{10}}_{\hand{10}} & X^{\deleted{11}}_{\hand{11}} & \ldots & X^{\deleted{1,n}}_{\hand{1n}}\\
\multicolumn{4}{l}{\ldots\ldots\ldots}
\end{array}
\]
\end{diagramblock}

formée des objets qui ligne par ligne et colonne par colonne satisfoht la
condition d'exctitude familière (transformer sommes amalgamées en produits
fibrés). De même la catégorie \xout{des} (Tricat) des (Bicat)-catégories
s'interprète en t ermes d'objets tri-\inserted[typed]{semi}simpliciaux, etc.

Exemple: soit C une catégorie, posons $C^{\deleted{o}}_{\hand{0}}=C$, $C^{\deleted{1}}_{\hand{1}}=\underline{Fl}(C)$, avec structures
$p_1,p_2$, évidentes (d'où une (Cat)-catétorie dont les composantes semi-sim-%
pliciales supérieures sont les catégories $\underline{Fl}^n(C)$ des n-flêches de C ...).
L'ensemble bi-semisimplicial correspondant a pour composantes les ensembles
Hom($\hand{\Delta}^n$x$\Delta^m$,C). Par exemple, les objets de $X^{\deleted{11}}_{\hand{11}}$ sont les carrés commutatifs
de flêches de C, qu'on peut interpréter(de deux manières différentes) comme

\mspage{\hand{36}}

\clearpage
% Tapis de Quillen, with Quillen's offprint and preprint (cote 162-5, typed part): archive cote 162-5, site page 37 (PDF page 38).
% typescript. Draft compared: none.
% AI-transcribed by Claude Opus 5.5 (claude-opus-5-5, adaptive thinking, effort high), Message Batches API, batch t162_full (msgbatch_01UAh3zyWPd454AeDkhnmQ33), 2026-11-13; not yet checked by eye.
\scanpage{162-5/37}
morphismes de flêches.

Notons une symmétrie évidente dans la c atégorie (Bicat), interprétée
en termes d\xout{\unsure{e}}'ensemble bi-semisimpliciaux, obtenue en échangeant les lignes
et les colonnes dans \xout{les} de tels ensembles bi... Un ensemble bisemisimplici-%
al qui ... peut donc de deux façons différentes être considéré comme défi-%
nissant une bicatégorie. De même le groupe symétrique $\mathrm{Sym}_n$ opère sur
la catégorie \xout{de} \unsure{(nupcat)|(n-upcat)} des catégories n-uples ..

Une 2-catégorie peut être définie comme constituée par un ensemble
\xout{discret} $C^{\hand{\unsure{0}}}_{\hand{0}}$ (l'ensemble des objets ), qu'on interprètera comme une catégo
rie discrète, plus la donnée pour tout x,y \hand{$\in$}$C^{\hand{0}}_{\hand{0}}$ d'une catégorie $\underline{\mathrm{Hom}}$(x,y),
plus \inserted[typed]{des foncteurs} \xout{des} composition entre celles\unsure{\hand{-}}ci ... Prenant la catégorie somme \xout{de} $C^{\hand{1}}_{\hand{1}}$
des $\underline{\mathrm{Hom}}$(x,y), On peut exprimer la donnée de ces foncteurs de composition,
et des sources et buts \inserted[typed]{x,y} dans les catégories $\underline{\mathrm{Hom}}$(x,y), comme définis par
des foncteurs ``source'' et ``but'' $C^{\hand{1}}_{\hand{1}} \hand{\to} C^{\hand{0}}_{\hand{0}}$, et un foncteur de composition
$C^{\hand{1}}_{\hand{1}} \mathrm{x}_{C^{\hand{0}}_{\hand{0}}} C^{\hand{1}}_{\hand{1}} \hand{\longrightarrow} C^{\hand{1}}_{\hand{1}}$ , les axiomes des 2-catégories correspondant exactement aux
axiomes des(Cat)-catégories. De cette façon, la catégorie des 2-catégories
est équivalente à la catégorie des (Cat)-catégories qui ont une ``catégorie
d'objets'' qui est discrète, ou encore à la catégorie des Bicatégories
qui, interprétées comme ensembles bi-semisimpliciaux, ont \xout{une} comme
première ligne un ensemble semisimplicial discret (i.e. un foncteur cons-%
tant sur \hand{$\Delta$} ). Il parait qu'on a une interprétation simple dans le même
esprit de la notion de n-catégorie en termes d'ensembles n-semi-simpliciaux
\margin{$\|$} \xout{xx} Demander à Quillen relations entre la notion de m-Catégorie dans
\margin{$\|$} la catégorie (n-Cat), et celle de ensemble (m+n)-semisimplicial.
\margin{$\|$} \xout{(en particulier entre cette dernière notion et celle de}

L'interprétation standard semisimpliciale permet d'interpréter

\mspage{\hand{37}}

\clearpage
% Tapis de Quillen, with Quillen's offprint and preprint (cote 162-5, typed part): archive cote 162-5, site page 38 (PDF page 39).
% typescript. Draft compared: none.
% AI-transcribed by Claude Opus 5.5 (claude-opus-5-5, adaptive thinking, effort high), Message Batches API, batch t162_full (msgbatch_01UAh3zyWPd454AeDkhnmQ33), 2026-11-13; not yet checked by eye.
\scanpage{162-5/38}
\margin{! Attention à la catégorie \uline{Hom} grab, qui s'\unsure{identitifie|identifie} à la catégorie \unsure{définie} par un $R\mathrm{Hom}_{Z}(C_{*},C_{*})$, ce point devrait être élucidé!}

de même les (Cat)-groupes(ou groupes à valeurs dans (Cat)) comme étant
les groupes semisimpliciaux qui "sont"des catégories i.e. qui \xout{sexdéfinisx}
\xout{sent} satsifont la condition d'exactitude familière, permettant d'en déter
miner tous les composants à l'aide de C$_{o}$\xout{$_{x}$} et C$_{1}$ . Appliquant le tapis de
Dold, dans le cas commutatif, on trouve que les (Cat)-groupes abéliens
forment une catégorie équivalente à celle des complexes de chaines de
groupes abéliens qui sont de longueur 1:
\[ 0 \leftarrow C_{o} \leftarrow C_{1} \leftarrow 0 \quad . \]
Quand on explicite \unsure{la|le} (Cat)-groupe abélien associé à un tel complexe de
chaines, on trouve "ma" construction chérie \unsure{da|de} la "catégorie de Picard"
associée au complexe: objets ceux de C$_{o}$ , flêches de x dans y les z dans C$_{1}$
tels que dz = y-x . On trouve donc à isomorphisme \unsure{ppès|près} tous les groupes
abéliens de (Cat) ainsi. Un complément important est que toute catégorie
\inserted[typed]{\deleted{(resp. et}}
de Picard \deleted{commutative)}est équivalente (au sens technique qu'on ne précise
\inserted[typed]{\deleted{(resp. (Cat)-groupe}}
pas ici) à \xout{laxcatégorie} un vrai (Cat)-groupe \deleted{(commutatif)} . \margin{Quand on}
\inserted{\unsure{homotopie} mod ``Gros'' équivalence, ce qui correspond à travailler dans la catégorie dérivée des complexes. Seuls invariants sont $H_0$, $H_1$ (\unsure{et l'invariant} $\pi_0\,\pi_1$).}
\margin{\unsure{Quitter définir} mais cela contredit les résultats de Verdier ? \illegible{3}}
\margin{\unsure{Mais si on prend} d'une cat. de Picard commutative, elle \unsure{ne vérifie pas} \unsure{(Cat)-groupe} \illegible{2}}
On notera qu'on a une \xout{lxxexp} classification beaucoup plus restreinte
que pour les H-espaces (resp. les H-espaces homotopiquement comm.) mod
équivalence d'homotopie, de façon précise entre les groupes resp. les grou-%
pes abéliens de la catégorie homotopique. La raison en est quon est beau\xout{xx}
\margin{//} coup plus restrictif dans le contexte catégorique pour le sens \unsure{da|de} l'associa
tivité (resp. et de la commutativité), puisqu'on ne demande pas seulement \margin{//}
l'existence d'isomorphismes d'associa\xout{i}tivité\xout{,} (qui suffirait pour avoir un
\inserted[typed]{la catégorie homotopique}
groupe dans \hand{$\gamma$}(Cat)' \hand{$\simeq$} (Ssimpl)', modulo l'existence d'unités et d'inverse)
mais qu'on suppose \uline{donné\unsure{\deleted{e}}s} ces isomorphismes, de façon à satisfaire à cer-%
taines compatibilités: c'est ce qui permet, me semble-t-il, de prouver qu'à
\margin{//} \margin{(} équivalence près, on peut supposer qu'il s'agit de vrais groupes dans
(Cat), \deleted{resp. de vraie commutativité.}

\mspage{\hand{38}}

\clearpage
% Tapis de Quillen, with Quillen's offprint and preprint (cote 162-5, typed part): archive cote 162-5, site page 39 (PDF page 40).
% typescript. Draft compared: none.
% AI-transcribed by Claude Opus 5.5 (claude-opus-5-5, adaptive thinking, effort high), Message Batches API, batch t162_full (msgbatch_01UAh3zyWPd454AeDkhnmQ33), 2026-11-13; not yet checked by eye.
\scanpage{162-5/39}
% NOTE: the right edge of the scan cuts the ends of several typed lines (e.g. "classifica", "qu'o", "semisimplici"); end-of-line hyphens may be hidden there.
On doit trouver une approximation meilleure d'un cran à la classifica
tion topologique en prenant des  (2-Cat)-groupes resp. des (2-Cat)-groupes
abéliens,  ou plus généralement des " Gr-2-catégories "  resp.
"G\unsure{ra}b-2-catégories" . (NB  U\unsure{ne} Grab-catégorie défi\xout{\unsure{n}}nit une Grab-2catégorie
comme une catégorie définit une 2-catégorie; du point de vue semi-simplici-%
al, cette opération serait analogue à celle de "cosquélette\xout{\unsure{"}} tronqué";
Mais également une Grab-2-catégorie définie en G\unsure{r}ab-catégorie\xout{\unsure{,}} savoir la
catégorie des automorphismes de l'objet unité; cette \xout{dé} opération serait
analogue à celle de "tuage des groupes d'homotopie en dim \hand{$>$}\over{\unsure{1}}{1}").  Pour la
notion gro\xout{\unsure{s}}ssière ( Gr\xout{\illegible{1}} ou Grab) l'idée de ma lettre à Verdier était qu'o
doit trouver la même chose que les tronqués\inserted{au cran 2}, dans la catégorie \uline{homotopique}
des H-espaces  -  i.e. ce qu'on déduit des H-espaces par tuage des groupes
d'homotopie en dimensions \hand{$>$}2 .

\margin{$\|$} (Demander à Quillen relations exactes entre \xout{H}objets-groupes et objets
\xout{\unsure{ponctués}}x connexes de la catégorie homotopique\inserted[typed]{ponctuée}: forment-ils des catégories équivalen-%
tes par le foncteur espace des lacets\hand{?}).
\margin{\unsure{non}, on a un foncteur \unsure{pl. fidèle} (?) \unsure{mais non ess.} \unsure{surjectif} ....}

Pour les \uline{vrais} groupes\inserted[typed]{abéliens} de (2-Cat), l'approche de Quillen consiste à
les identifier \xout{\unsure{à}}d'abord à des groupes abéliens de la catégorie (Bicat),
donc à des groupes abéliens bi-semi\xout{e}simpliciaux particuliers (conditions
sur chaque ligne et chaque colonne). Appliquant la version bi-semisimplici
ale du tapis de Dold, on trouve que \xout{\illegible{1}} la catégorie de ces objets
estb équivalente à la catégorie des \xout{objets} bicomplexes \inserted[typed]{de chaines} de groupes  abélien
qui sont nuls en dehors du carré  \unsure{0}\hand{$\leq$}i,j \hand{$\leq$} 1.  \xout{Quillen dit que} La colonne
que la première ligne corresponde à un \xout{\unsure{ensemble}} groupe semisimplicial
trivial signifie alors qu'elle n'a que le terme de degré zéro, i.e. que
le bicomplexe est réduit à

\begin{diagramblock}\begin{tikzcd}
C_{oo} & 0 \arrow[l] & \hand{$0$} \arrow[l] \\
C_{1o} \arrow[u] & C_{11} \arrow[l] \arrow[u] & \hand{$0$} \arrow[l]
\end{tikzcd}\end{diagramblock}
% NOTE: the bottom edge of the scan cuts just below the diagram; a mark under C_{1o} is partly hidden.

\mspage{39}

\clearpage
% Tapis de Quillen, with Quillen's offprint and preprint (cote 162-5, typed part): archive cote 162-5, site page 40 (PDF page 41).
% typescript. Draft compared: none.
% AI-transcribed by Claude Opus 5.5 (claude-opus-5-5, adaptive thinking, effort high), Message Batches API, batch t162_full (msgbatch_01UAh3zyWPd454AeDkhnmQ33), 2026-11-13; not yet checked by eye.
\scanpage{162-5/40}
Ces\xout{x} bicomplexes s'identifient aux complexes de chaines de longueur 2.
Donc la catégorie de (2-Cat)-groupes abéliens est équivelente à celle des
complexes de chaines de longueur 2. L'explicitation de la 2-catégorie
associée à un tel complexe de chaines est mon interprétation chérie. Il
faudrait encore expliciter la\xout{sxmorphismex} 2-Grab-catégorie des
homomorphismes (en un sens à préciser) entre deux\xout{x} telles 2-Grab-catégo-%
ries strictes, en termes d'un complexe RHom.

L'idée qui se dégage de cette approche, et que Quillen on espère va
expliciter, c'est que la \xout{notionxdexcatégoriexde} notion\xout{s} de groupe abélien
dans (n-Cat) correspond à celle de complexe de \inserted{chaines de} longueur n, en raisonnant
soit mod. isomorphisme de part et d'autre, soit mod. une notion naturelle
de n-équivalence d'un côté, et dans la catégorie dérivée de l'autre. Ceci
doit pouvoir se préciser en utilisant des Hom$^{\bullet}$ resp. des RHom$^{\bullet}$.

Si maintenant C est une catégorie avec un ``produit tensoriel'' associa
tif (et éventuellement unitaire), i.e. un bifoncteur avec des isomorphis-%
mes d'associativité (mais sans inverse ...), on peut lui associer sans
doute une Gr-catégorie universelle (tout comme à un monoïde on associe un
groupe symmétrisé), d'où des invariants, qu'on pourra désigner par $K_o(C)$
et $K_1(C)$. On suppose qu'on pourra de même trouver, pour tout n, une
Gr-n-catégorie universelle où envoyer C de façon compatible avec $\otimes$ (tout
cela dans un sens à préciser), dont les invariants $\hand{$\pi$}_n$ seraient par défini
tion les $K_n(C)$. Un cas intéressant est celui où C est une catégorie avec
produits finis (par exemple une catégorie additive) et $\otimes$ est le produit,
\xout{maisx} \over{(}{[}Mais on notera que dans le cas d'une catégorie non seulement additive
mais abélienne, il y a une notion plus adéquate du $K_o$ , qui devrait corres
pondre à une notion plus adéquate de la suite des $K_i$ , inspirée du tapis
de Roos; sans doute c'est en termes de (``vraie'') catégories triangulées

\mspage{40}

\clearpage
% Tapis de Quillen, with Quillen's offprint and preprint (cote 162-5, typed part): archive cote 162-5, site page 41 (PDF page 42).
% typescript. Draft compared: none.
% AI-transcribed by Claude Opus 5.5 (claude-opus-5-5, adaptive thinking, effort high), Message Batches API, batch pilotnext (msgbatch_01A9sdu8SYq23bEtqX9tUncQ), 2026-11-13; not yet checked by eye.
\scanpage{162-5/41}
qu'il conviendra de poser les définitions adéquates). - Peut-être Quillen
\margin{oui} \margin{$\|$}a-t-il une définition semi-simpliciale plus directe, sans passer par des
n-catégories, des $K_i$ et plus précisément du H-espace qui devrait correspon-%
dre à C. \xout{Dans le cas} Ce n'est certainement pas \inserted[typed]{en général} l'ensemble semi-simplicial
SC \xout{engendré à} associé à C, puisque dans le cas d'une catégorie additive,
C est homotopiquement trivial (en effet le foncteur identique est ``homotope
au foncteur nul, puisqu'il suffit d'avoir un homomorphisme d'un foncteur
dans l'autre; d'où une équivalence d'homotopie entre C et la catégorie
nulle). Pourtant, il semble bien que dans le cas où C est une Gr-catégorie
\margin{\unsure{$\xi$|$\zeta$}} \margin{$|$}i.e. admet des inverses, alors ce soit bien la construction ``naïve'' SC qui
est la bonne, puisqu'on trouve bien un $\hand{\pi}_0$ et un $\hand{\pi}_1$ qui sont aussi les
$\hand{\pi}_0$ et $\hand{\pi}_1$ au sens  catégorie sans Gr-structure.

\margin{À expliquer : pour des Gr-n-catégories, les $\pi_i$ au sens
\unsure{$\xi$|$\zeta$} $|$ semi-simplicial sont aussi les $\pi_i$ au sens Gr-n-catégories,
savoir les objets \unsure{:|;} isom pris dans divers $\underline{\mathrm{Hom}}(1,1)$ ...... \deleted{Il y a
aussi une interprétation de l'opération \unsure{$\Omega$ : espace des lacets}
en termes d'opération \unsure{de|du} \unsure{shift} géométrique évidente sur les
n-catégories .....}}

\mspage{41}

\clearpage
% Tapis de Quillen, with Quillen's offprint and preprint (cote 162-5, typed part): archive cote 162-5, site page 42 (PDF page 43).
% typescript. Draft compared: none.
% AI-transcribed by Claude Opus 5.5 (claude-opus-5-5, adaptive thinking, effort high), Message Batches API, batch t162_full (msgbatch_01UAh3zyWPd454AeDkhnmQ33), 2026-11-13; not yet checked by eye.
\scanpage{162-5/42}
% NOTE: no typed text and no other writing on this page; it is blank apart from the number at the foot
\mspage{42}

\clearpage
% Tapis de Quillen, with Quillen's offprint and preprint (cote 162-5, typed part): archive cote 162-5, site page 43 (PDF page 44).
% typescript. Draft compared: none.
% AI-transcribed by Claude Opus 5.5 (claude-opus-5-5, adaptive thinking, effort high), Message Batches API, batch t162_full (msgbatch_01UAh3zyWPd454AeDkhnmQ33), 2026-11-13; not yet checked by eye.
\scanpage{162-5/43}
III   Point de vue "motivique" en théorie du cobordisme.

Soit C la catégorie des variétés différentiables (pas nécessairement
orientables), \inserted[typed]{les morphismes étant les classes d'homotopie d'applications continues.} Si B est une catégorie, on s'intéresse aux couples
$(F_{.},F^{\cdot})$ d'un foncteur covariant et d'un foncteur contravariant de C
dans B, satisfaisant les conditions que pour tout X \hand{$\in$} Ob C, on a $F_{.}(X)=$
$F^{\cdot}(X)$, et que si on a un produit fibré ordinaire
\begin{diagramblock}
\begin{tikzcd}[column sep=small, row sep=small]
 & X' \arrow[dl, "g'"'] \arrow[dr, "f'"] & \\
X \arrow[dr, "f"'] & & Y' \arrow[dl, "g"] \\
 & Y &
\end{tikzcd}
\end{diagramblock}
,
avec f et g transversaux, alors on a $g^{\cdot}f_{.} = f!\ g'^{\cdot}$ \inserted[typed]{(on pose $f_{.}=F_{.}(f),f^{\cdot}=F^{\cdot}(f)$)}. La question est
de trouver une B et un couple $(F_{.},F^{\cdot})$ qui soient "universels" dans un
sens évident. Quillen prouve que la construction suivante donne une
telle solution.

Pour toute variété X\hand{$\in$} Ob C , on définit un anneau commutatif
B(X) (l'anneau de bordisme de X) en prenant comme éléments les classes
\xout{de} d'éléments de C au dessus de X, pour la relation de cobordisme:
$f':X'\to X$ et $f'':X''\to X$ sont équivalents, s'il existe une variété à bord
Z, et un $f:Z\to X$, tel que bord $Z = X'$\hand{$\amalg$}$X''$ et f' et f" sont induits
par f. La somme de B(X) est la somme disjointe (c'est une loi de groupe,
\xout{l'inverse étant} où tout élément est d'ordre \deleted{\unsure{1}}\hand{2} \deleted{\unsure{ou 2}}), le produit est
donné par le produit fibré, en utilisant le théorème de trahsversalité
de Thom pour bouger un peu f", de sorte à le rendre transversal à f'.
\inserted{\uline{NB} L'anneau est de plus \unsure{gradué} par la dimension relative ...}
Si on a $f: X\to Y$, on définit $f_{.} : B(X)\to B(Y)$ de façon évidente (toute
variété sur X est sur Y), et $f^{\cdot}: B(Y)\to B(X)$ par produit fibré (utili-%
sant encore le th. de transversalité). De cette façon, B(X) devient un
foncteur-anneau \inserted{(gradué)} contravariant en X, et un foncteur groupe covariant en

\mspage{43}

\clearpage
% Tapis de Quillen, with Quillen's offprint and preprint (cote 162-5, typed part): archive cote 162-5, site page 44 (PDF page 45).
% typescript. Draft compared: none.
% AI-transcribed by Claude Opus 5.5 (claude-opus-5-5, adaptive thinking, effort high), Message Batches API, batch t162_full (msgbatch_01UAh3zyWPd454AeDkhnmQ33), 2026-11-13; not yet checked by eye.
\scanpage{162-5/44}
X. De plus, on a la formule \inserted{(du type $g'f. = f'.g'$ ci-dessus, impliquant la formule} de projection habituelle. Ceci posé, les
constructions bien connues en théorie des correspondances permettent
de définir une catégorie $\underline{B}$ dont \unsure{lès|les} objets sont ceux de C, et les
Hom étant donnés par
\[ \mathrm{Hom}_{B}(X,Y) = B(XxY) \quad . \]
On a un isomorphisme canonique $\underline{B} \simeq \underline{B}^{o}$ , qui est l'identité sur les
objets et la symmétrie $B(XxY)\simeq B(YxX)$ sur les flêches, et un foncteur
évident \xout{Bxxx$\underline{B}$x} $F. : C \to \underline{B}$ , qui est l'identité sur les objets et donné
par les graphes sur les flêches. Composé avec la symmétrie précédente,
on trouve un foncteur $F^{\cdot}$: \xout{BX} $C$ \_ $\underline{B}$. On vérifie enfin que le couple
$(F.,F^{\cdot})$ satisfait aux conditions envisagées dans l'alinéa 1 (ceci aussi
devrait être formalisé au cas des correspondances générales ...). \xout{On}
vérifie enfin qu'il a la propriété universelle qu'il faut. Lorsqu'on
se donne un couple $(F!,F'^{\cdot})$ à val eursdans une catégorie B', il faut
trouver une factorisation par B. Elle est toute trouvée sur les
objets, il faut la définir sur les flêches, donc définir, pour X,Y
dans C, une application
\[ (*) \qquad B(XxY) \to \mathrm{Hom}(X',Y') \ . \]
Mais si Z est au\unsure{d}dessus de XxY, d'où $p:Z\to X$ et $q:Z\to Y$, on lui associe
un élément du deuxième membre de (*) comme le composé
\[ X' \overset{F'^{\cdot}(p)}{\hand{$\longrightarrow$}} Z' \overset{F!(q)}{\hand{$\longrightarrow$}} Y' \quad . \]
Il faut prouver que l'élément obtenu en \unsure{remplaçant} Z \inserted{$=Z_o$} par un $Z_1$ cobor\xout{d}
dant au dessus de X \inserted{est le même}. Or si $g: T\to XxY$ réalise le cobordisme, on
voit (regardant seulement T et la décomposition bord(T) = $Z_{o}\amalg Z_1$)
qu'il existe une variété $\overline{T}$ sur $S^1$ (obtenue en ``doublant'' T) \xout{xtxxxx}
\xout{applixatixxxTxxxxx} transversale au dessus d'un couple de poihts $s_o,s_1$
de $S^1$, tels que $Z_i$ soit l'image inverse \inserted{dans $\overline{T}$} de $s_i$ (i=0,1); d'autre part,

\mspage{44}

\clearpage
% Tapis de Quillen, with Quillen's offprint and preprint (cote 162-5, typed part): archive cote 162-5, site page 45 (PDF page 46).
% typescript. Draft compared: none.
% AI-transcribed by Claude Opus 5.5 (claude-opus-5-5, adaptive thinking, effort high), Message Batches API, batch t162_full (msgbatch_01UAh3zyWPd454AeDkhnmQ33), 2026-11-13; not yet checked by eye.
\scanpage{162-5/45}
g manifestement se prolonge en $\overline{g} : \overline{T} \to X\mathrm{x}Y$ . Cela nous ramène à prouver
\inserted[typed]{le composé}
que \hand{\unsure{$\Gamma$}}$\overline{T}'$\hand{$\xrightarrow{\alpha_i}$}\hand{\unsure{$\Gamma$}}$Z_i^!$\hand{$\xrightarrow{\alpha_{i*}}$}$\overline{T}'$ ne dépend pas de i=0,1. Ecrivons X au lieu de $\overline{T}'$,et
considérons le diagramme cartésien d'inclusions
\begin{diagramblock}\begin{tikzcd}
X_s \arrow[r, "\hand{\alpha_s}"] \arrow[d, "\hand{\alpha_s}"'] & X \arrow[d, "\hand{\Gamma = (p,\mathrm{id}_X)}"] \\
X \arrow[r, "i_s"] & S^1\mathrm{x}X
\end{tikzcd}\end{diagramblock}
,

on a alors $\hand{\alpha}^{\hand{*}}_s \hand{\alpha}_{s*} = \Gamma^* i_{s\hand{*}}$ , qui ne dépend pas de s puisque $i_s$ n'en
dépend pas à homotopie près. - Le reste est sans doute l'AQT.

Sur la strcture des B.(X). Ce sont tous des algèbres graduées
sur B.(pt)=\hand{$\Lambda$} , qui est une algèbre graduée à degrés positifs, réduite
\margin{\unsure{\S}}
à $\underline{F}_2$ en degré\xout{s} 0. C'est d'ailleurs une algèbre de polynômes \inserted{(à une infinité de générateurs \unsure{$x_i$ ($x_i$ de degré de la forme $2^i-1$)}} \deleted{\unsure{connue}} ...
\inserted[typed]{d'\unsure{anneaux gradués}}
On a un homomorphisme évident B.(X) $\to$ H$^{\over{\cdot}{\unsure{*}}}$(X), nul en degrés \hand{$>$}0, qui
à tout Z sur X associe l'image de 1\hand{$\in$}H$^0$(Z) par l'homomorphisme de Gysin
associé à Z$\to$ X. Cet homomorphisme par Thom est surjectif (toute clas-%
\margin{en fait $H(X) \simeq B.(X)/\Lambda^+B.(X)$}
se de cohomologie provient d'une image de variété compacte ...); \xout{noter}
On ppouve même qu'il existe un relèvement fonctoriel en X
\margin{?}
H$^\bullet$(X) $\to$ B.(X), \inserted{il est douteux qu'il soit}\hand{(}compatible avec multiplication\hand{)} \xout{tel}tel que l'on ait\hand{:}

\over{a)}{\unsure{1)}} L'homomorphisme correspondant H$^\bullet$(X)$\otimes_{\underline{F}_2}$\hand{$\Lambda$} \hand{$\simeq$}B.(X) est un iso-%
\margin{?}
morphisme (ce sera vrai pour tout relèvement individuellement ...), \deleted{et}

\over{b)}{\unsure{2)}} \hand{[}\inserted{Quillen ignore si} \over{Le\xout{x}}{\unsure{le}} relèvement est compatible aussi avec les hom. de Gysin.\hand{]}

Ainsi, la catégorie \underline{B} admet un foncteur pleinement fidèle dans la
catégorie des modules gradués libres de type fini sur \hand{$\Lambda$}, les modules
de type fini qu'on obtient ainsi ayant des générateurs à degrés \hand{$<$} 0 et
satisfaisant (pour X connexe i.e. quand il y a un seul générateur de
degré zéro\hand{)} la dualité de Poincaré. Quand on prend la catég\unsure{o|b}rie quasi-%
abélienne\hand{-}\inserted{\unsure{B}}associée à \underline{B}, on trouve donc tous les modules gradués libres
de type fini à générateurs \xout{négatifs} en degrés \hand{$\leq$} 0 seulement. C'est une

\mspage{45}

\clearpage
% Tapis de Quillen, with Quillen's offprint and preprint (cote 162-5, typed part): archive cote 162-5, site page 46 (PDF page 47).
% typescript. Draft compared: none.
% AI-transcribed by Claude Opus 5.5 (claude-opus-5-5, adaptive thinking, effort high), Message Batches API, batch t162_full (msgbatch_01UAh3zyWPd454AeDkhnmQ33), 2026-11-13; not yet checked by eye.
\scanpage{162-5/46}
catégorie \inserted[typed]{additive(et même quasi-abélienne)}avec $\otimes$ qui a toutes les vertus,sauf d'être abélienne, et
sauf d'admettre des \uline{Hom} internes. Pour trouver ces derniers, on peut
si on veut paraphraser la construction de la catégorie des motifs, en
\xout{\unsure{introduisant un motif}} considérant l'objet \inserted{$\Lambda(-1)$} de $\tilde{\underline{B}}$ conoyau de
$p^{\hand{*}}$ : $e \hand{$\longrightarrow$} S$, où $p:S^{1} \hand{$\longrightarrow$} e$ est l'application du cercle sur un
point. C'est un objet de poids 1 (le poids est défini par le degré \xout{d}
des générateurs, changés de signe). On constate aussitôt que le
foncteur $M\hand{$\mapsto$} M(-1)=M\otimes\hand{$\Lambda$}(-1)$ de $\tilde{\underline{B}}$ dans lui-même est pleinement fidèle,
ce qui permet alors, par une construction formelle standard, de complé-%
ter $\overset{\hand{$\sim$}}{\underline{B}}$ en une catégorie où $\hand{$\Lambda$}(-1)$ devient inversible. Il se trouve
alors que dans cette dernière les \uline{Hom}(X,Y) internes existent et
\xout{satisfont} sont isomorphes à $\check{X}\otimes Y$ ... en fait , on trouve maintenant la
catégorie de\xout{s} \uline{tous} les modules gradués libres de type fini sur \quad .
Bien sûr, cette catégorie n'est pas encore abélienne pour autant,
puisque \quad n'est pas un corps; \xout{pour} \inserted[typed]{si on veut} définir la catégorie abélienne
engendrée, on s'attend à trouver la\xout{s} catégorie \hand{/}desmodules de \deleted{type} \inserted{présentation} fin-%
sur $\hand{$\Lambda$}$.( \inserted{\unsure{qui est abélienne bien que} $\Lambda$ \unsure{ne soit pas noethérien, car} $\Lambda$ \unsure{est}} \deleted{\quad est-il noethérien, i.e. le nombre de générateurs est-il}
\margin{\unsure{?}} \deleted{fini ? ).} \hand{\unsure{limite d'anneaux noeth. à homs de transition} \uline{plats})}.

Cette construction ressemblerait à celle d'une théorie des
motifs où on n'aurait pas passé à l'équivalence homologique pour
les cycles, mais où on aurait gardé l'équivalence rationelle, - qui
est bien trop fine \xout{\unsure{pou}}à beaucoup d'égards. Pour trouver un équiva\xout{l}%
lent plus étroit, il convient de reprendre la théorie prédédente,on
pourrait formuler le même Pb universel que plus haut, mais en prenant
l'équivalence homologique (via le graphe) des morphismes de variétés,
au lieu de l'équivalence homotopique. \deleted{Mais alors il semble que l'on
trouve comme catégorie solution celle dont les objets sont les varié-}

\mspage{46}

\clearpage
% Tapis de Quillen, with Quillen's offprint and preprint (cote 162-5, typed part): archive cote 162-5, site page 47 (PDF page 48).
% typescript. Draft compared: none.
% AI-transcribed by Claude Opus 5.5 (claude-opus-5-5, adaptive thinking, effort high), Message Batches API, batch t162_full (msgbatch_01UAh3zyWPd454AeDkhnmQ33), 2026-11-13; not yet checked by eye.
\scanpage{162-5/47}
\deleted{tés, les morphismes \uline{toutes} les correspondances \inserted[typed]{co}homologiques (en vert\unsure{u}}
\deleted{du thorème de Thom rappelé plus haut). On trouve donc un foncteur}
\deleted{\unsure{pleinemeht|pleinement} fidèle de cette c atégorie dans celle d es vectoriels}
\deleted{gradués \xout{\unsure{dexd}} positivement de dim. finie sur $\underline{F}_2$ , qui est déjà abélien}
\deleted{ne, et si on veut rendre comme plus haut $\underline{F}_2(-1)$ inversible, on trouve}
\deleted{tous les vectoriels gradués de dim. finie sur $\underline{F}_2$ . C'est un peu trop}
\deleted{trivial comme structure, et on peut faire mieux, \xout{en prenant comme}}
\inserted[typed]{Pour trouver une catégorie universelle,}
\xout{\unsure{correspondances. Pour ceci}} on associe à toute X \xout{\unsure{le sous-anneau ensemble}}
\xout{\unsure{de}} l'ensemble \inserted[typed]{S(X)} des sous-variétés \xout{S(X)} de X, modulo équivalence homolo%
gique. Donc S(X) \hand{$\subset$} $H^{\cdot}(X)$, et on prouve qu'on trouve en fait un sous-%
anneau de $H^{\cdot}(X)$ (grâce à un théorème de Thom qui caractérise S(X) en
termes d'opérations de Steenrod, cf plus bas). On prouve de même
(par le même théorème) que les S(XxY) sont stables par composition de\unsure{s}
correspondances cohomologiques( cela ne semble pas trivial, puisque
les S(X) ne sont pas stables par Gysin...). D'oû une catégorie addi-%
tive dont les objets sont ceux de C, les flêches sont les S(XxY).
Il semble qu'elle doive jouer le rôle universel que j'ai envisagé
plus haut. En tous cas, sa définition géom\xout{m}étrique est fort simple !
Nous allons en donner une interpétation algébrique. Pour ce\unsure{ci},
considérons l'anneau des opérations cohomologiques \inserted[typed]{``stables''} \deleted{\xout{additives sur}} \deleted{i.e.}
\margin{vertical line in the left margin, from this line down to the end of the typed text}\deleted{\inserted[typed]{des endomorphismes du}}
\margin{\unsure{$\xi$}}\deleted{\xout{le} foncteur X \quad $H^{\cdot}$(X) qui sont additifs, et qui en degré zéro se}
\deleted{réduisent à un multiple de l'identité.} C'est un anneau gradué, avec
application diagonale anticommutative, \inserted[typed]{i.e. commutative (car. 2 !)} associative et unitaire. Son
dual gradué est une algèbre \inserted[typed]{graduée anti\unsure{c}commutative donc \xout{commutative},} commutative,avec application diagonale
associative et unitaire; il est réduit à $\underline{F}_2$ en degré zéro. Donc il
corespond à un schéma en groupes \unsure{aff}ine G sur $\underline{F}_2$ , qui est d'ailleurs
le foncteur des \unsure{$\otimes$}-automorphismes \inserted[typed]{additifs} du foncteur $H^{\cdot}$ sur C \inserted{qui \unsure{opèrent}
trivialement sur \deleted{\unsure{l'objet}} $H^1(S^1)$. \deleted{\illegible{1}}}

\mspage{47}

\clearpage
% Tapis de Quillen, with Quillen's offprint and preprint (cote 162-5, typed part): archive cote 162-5, site page 48 (PDF page 49).
% typescript. Draft compared: none.
% AI-transcribed by Claude Opus 5.5 (claude-opus-5-5, adaptive thinking, effort high), Message Batches API, batch pilotnext (msgbatch_01A9sdu8SYq23bEtqX9tUncQ), 2026-11-13; not yet checked by eye.
\scanpage{162-5/48}
Le théorème de Thom est alors qu'un $x \hand{$\in$} H^{\cdot}(X)$ est dans S(X) sss il est
\xout{xxxxxx}\inserted[typed]{fixe.} par G, i.e. satisfait $A^{+}. x = 0$. (Pourquoi cette condition est-el% NOTE: right edge of the page cut, end of line not visible
le déjà nécessaire ?) Une autre façon de le voir est de dire que le
foncteur $H^{\cdot}$ considéré comme allant de la catégorie $\underline{S}$ précédente dans
la catégorie des \xout{\unsure{vectoriels}} représentations de G dans des vectoriels
sur $\underline{F}_2$ , est pleinement fidèle. \xout{\unsure{Comme G est}} Quand on prend l'enveloppe
pseudo-abélienne, on trouve une sous-catégorie pleine de la catégorie des
représentations de G, mais il est douteux qu'on trouve une catégorie abé-%
lienne, voire la catégorie de toutes les représentations de G de dimensio% NOTE: right edge cut
finie sur $\underline{F}_2$ . (Quelle est la structure de G, qui est prolisse; d'après
\margin{oui, car \unsure{ga-} \illegible{3} \illegible{3} \illegible{2} \deleted{\illegible{1}} \illegible{4} \illegible{3} \unsure{fini stable par} $\triangleright$}Quillen, il serait pro-unipotent ....). Ici G bien sûr joue le rôle d'un
groupe de Galois motivique, mais contrairement à la situation familière,
il semble qu'il soit fortement non réductif (et d'ailleurs le corps de
définition dans tout ceci est $\underline{F}_2$ , et non $\underline{Q}$ ! ).

Quillent donne une généralisation de la construction d'une caté-%
gorie universelle $\underline{B}$ au cas où on travaille avec des variétés pas nécessai% NOTE: right edge cut
rement propres (mais f. n'étant supposé défini que pour f propre), et
quand on introduit une condition d'orientation \inserted[typed]{stable} sur les \uline{morphismes}
envisagés , via un morphisme $H \hand{$\to$} BO(\infty )$, où H est un H-espace (commuta-%
tif, associatif et unitaire ?): une orientation de $f:X \hand{$\to$} Y$ sera un
relèvement homotopique de $X \hand{$\to$} BO$ défini par \xout{fxx(Tx} $T_X - f^{\cdot}(T_Y)$ en $X \hand{$\to$} H$.
les éléments de H(X,Y) seront définis par des diagrammes de morphismes
orientés $Z \hand{$\to$} X$, $Z \hand{$\leftarrow$} Y$, le deuxième étant propre, modulo bords de tels.
En tant qu'ensemble, il se définit aussi en termes de la vari té XxY
munie simplement de la famille des supports formée des ensembles fermés
propres sur Y ..... \hand{Les $B^{H}_{m}(X)$ (X compact disons) peuvent}
\hand{se décrire homotopiquement à la Thom, comme les groupes d'homotopie}
\hand{stables d'espaces \uline{Hom}(X, \unsure{$MH(N-m)$})} \hand{-----.}

\mspage{\hand{48}}

\clearpage
% Tapis de Quillen, with Quillen's offprint and preprint (cote 162-5, typed part): archive cote 162-5, site page 50 (PDF page 51).
% offprint. Draft compared: none.
% AI-transcribed by Claude Opus 5.5 (claude-opus-5-5, adaptive thinking, effort high), Message Batches API, batch t162_full (msgbatch_01UAh3zyWPd454AeDkhnmQ33), 2026-11-13; not yet checked by eye.
\scanpage{162-5/50}
Inventiones math. 5, 237 -- 242 (1968)

On the Representation of Hermitian Forms
as Sums of Squares

Daniel G. Quillen* (Cambridge, Mass.)

There is an extensive literature (e.g. [1], Ch.VI) concerned with
Hilbert's problem of representing a positive-valued form with real
coefficients as a sum of squares. The corresponding problem for her-%
mitian forms does not seem to have been touched. In this paper we
present a theorem about hermitian forms which arose in connection
with a priori estimates for partial differential equations. This theorem is
proved by the methods of a priori estimates and one of its corollaries
is the Hilbert Nullstellensatz over $\mathbf{C}$, so we obtain a direct analytic proof
of the latter theorem.

Let $T$ and $V$ be finite dimensional complex vector spaces and let
$E(u_1,u_2)$ be a hermitian form on $S_m T\otimes V$ where $S_m$ denotes $m$-th sym-%
metric tensors. If $\zeta\in T$ and $v\in V$ let
\[ E^{\wedge}(\zeta,v)=(m!)^{-2}\,E(\zeta^m\otimes v,\zeta^m\otimes v). \]
Then $E^{\wedge}$ is a real-valued smooth function on $T\times V$ such that
\[ E^{\wedge}(s\,\zeta,t\,v)=|s|^{2m}\,|t|^2\,E^{\wedge}(\zeta,v) \qquad (1) \]
for $s,t\in\mathbf{C}$, and the map $E\mapsto E^{\wedge}$ yields a 1-1 correspondence between
hermitian forms on $S_m T\otimes V$and smooth functions on $T\times V$satisfying (1).
The inverse procedure of passing from a function to a form is polariza-%
tion (see formulas (2) below).

Let $|\zeta|^2$ be a hermitian scalar product on $T$ and define the ``$N$-th
prolongation'' of $E$ to be the hermitian form $E^N$ on $S_{m+N}T\otimes V$ with
\[ (E^N)^{\wedge}(\zeta,v)=(N!)^{-1}\,|\zeta|^{2N}\,E^{\wedge}(\zeta,v). \]
Theorem. The following conditions are equivalent:
(i) $E^{\wedge}(\zeta,v)>0$ for $\zeta\neq 0$ and $v\neq 0$.
(ii) For sufficiently large $N$, $E^N$ is positive definite.

We first derive an expression for $E^N$ as an integral. By choosing an
orthonormal basis for $T$ we may identify $T$ with the space of $\mathbf{C}$-linear
functions on $\mathbf{C}^n$ in such a way that the orthonormal basis is $z^1,\dots,z^n$
where $z^i$ is the $i$-th coordinate function. Then $S_k T$ may be naturally
identified with the holomorphic polynomial functions homogeneous of

* This research was supported in part by NSF GP-6959.

17 Inventiones math., Vol. 5

\mspage{\hand{50}}

\clearpage
% Tapis de Quillen, with Quillen's offprint and preprint (cote 162-5, typed part): archive cote 162-5, site page 51 (PDF page 52).
% offprint. Draft compared: none.
% AI-transcribed by Claude Opus 5.5 (claude-opus-5-5, adaptive thinking, effort high), Message Batches API, batch t162_full (msgbatch_01UAh3zyWPd454AeDkhnmQ33), 2026-11-13; not yet checked by eye.
\scanpage{162-5/51}
\mspage{238} \hfill D.G. Quillen:

degree $k$ and $S_k T$ has the basis $z^\alpha$, $|\alpha|=k$ where $\alpha$ is a multi-index. If we
identify $V$ with $\mathbf{C}^r$ and let $e_1,\dots,e_r$ be the canonical basis of $\mathbf{C}^r$, then
$S_k T\otimes V$ can be naturally identified with the holomorphic polynomial
functions homogeneous of degree $k$ with values in $\mathbf{C}^r$. Then $S_m T\otimes V$
has the basis $z^\alpha e_i$, $|\alpha|=m$, $1\leqq i\leqq r$, so there are complex numbers $a_{\alpha i,\beta j}$,
$|\alpha|=|\beta|=m$, $1\leqq i, j\leqq r$ with $a_{\beta j,\alpha i}=a^*_{\alpha i,\beta j}$ such that
\[
\begin{aligned}
E(u_1,u_2)&=\sum \alpha!\,\beta!\,a_{\alpha i,\beta j}\,b_{\alpha i}\,c^*_{\beta j}\\
E^\wedge(\zeta,v)&=\sum a_{\alpha i,\beta j}\,\zeta^\alpha\,v_i\,\zeta^{*\beta}\,v^*_j
\end{aligned}
\tag{2}
\]
if $u_1=\sum b_{\alpha i}\,z^\alpha e_i$, $u_2=\sum c_{\alpha i}\,z^\alpha e_i$, $\zeta=\sum \zeta^i z^i$, and $v=(v_1,\dots,v_r)$, where *
denotes complex conjugate. The matrix $a_{\alpha i,\beta j}$ may be determined from
the Taylor series of $E^\wedge$ at $(0,0)$ which shows that $E$ and $E^\wedge$ determine
each other.

Let $E^\sim$ be the differential quadratic form
\[
E^\sim(f,g)=\sum a_{\alpha i,\beta j}(D^\alpha f_i)(D^\beta g_j)^*
\]
where $f$ and $g$ are holomorphic functions on $\mathbf{C}^n$ with values in $\mathbf{C}^r$ and
$D^i=\partial/\partial z^i$. If $dG=\pi^{-n}\exp(-|z|^2)\,dx^1\,dy^1\dots dx^n\,dy^n$ is the Gaussian
measure on $\mathbf{C}^n$, then we have the following expression for $E^N$.

Lemma 1. If $f,g\in S_{m+N}T\otimes V$, then
\[
E^N(f,g)=\int E^\sim(f,g)\,dG.
\]

Proof. If $u_1$, $u_2$ are two holomorphic polynomial functions on $\mathbf{C}^n$, let
\[
(u_1,u_2)_0=\int u_1\,u_2^*\,dG.
\]

By integrating by parts we obtain the basic formula
\[
(D^i u_1,u_2)_0=(u_1,z^i u_2)_0,
\]
from which follows
\[
(z^\alpha,z^\beta)_0=\begin{cases}0 & \alpha\neq\beta\\ \alpha! & \alpha=\beta.\end{cases}
\tag{3}
\]
By definition
\[
\begin{aligned}
(E^N)^\wedge(\zeta,v)&=(N!)^{-1}\,|\zeta|^{2N}\,E(\zeta,v)=\sum_{|\gamma|=N}(\gamma!)^{-1}\,\zeta^\gamma\,\zeta^{*\gamma}\,E(\zeta,v)\\
&=\sum(\gamma!)^{-1}\,a_{\alpha i,\beta j}\,\zeta^{\gamma+\alpha}\,v_i(\zeta^*)^{\gamma+\beta}\,v_j^*,
\end{aligned}
\]
hence
\[
(E^N)^\sim(f,g)=\sum_{|\gamma|=N}(\gamma!)^{-1}\,E^\sim(D^\gamma f,D^\gamma g).
\]
By linearity it suffices to prove the lemma when $f=z^\delta e_i$ and $g=z^\varepsilon e_j$
where $|\delta|=|\varepsilon|=m+N$. But
\[
E^N(f,g)=(E^N)^\sim(f,g)(0)=\sum_{\substack{|\alpha|=|\beta|=m\\ |\gamma|=N}}(\gamma!)^{-1}\,a_{\alpha i,\beta j}(D^{\gamma+\alpha}z^\delta)(D^{\gamma+\beta}z^\varepsilon)^*
\]

\mspage{\hand{51}}

\clearpage
% Tapis de Quillen, with Quillen's offprint and preprint (cote 162-5, typed part): archive cote 162-5, site page 52 (PDF page 53).
% offprint. Draft compared: none.
% AI-transcribed by Claude Opus 5.5 (claude-opus-5-5, adaptive thinking, effort high), Message Batches API, batch t162_full (msgbatch_01UAh3zyWPd454AeDkhnmQ33), 2026-11-13; not yet checked by eye.
\scanpage{162-5/52}
On the Representation of Hermitian Forms as Sums of Squares \mspage{239}

and
\[ \int E^{\sim}(f,g)\,dG = \sum_{|\sigma|=|\tau|=m} a_{\sigma i,\tau j}(D^{\sigma} z^{\delta}, D^{\tau} z^{\varepsilon})_0 . \]
These two coincide by (3), proving the lemma.

Introduce norms
\[ \|u\|_q^2 = \int \sum_{|\gamma|=q} (\gamma!)^{-1}\, |D^{\gamma} u|^2\, dG \]
\[ \|u\|_{q*}^2 = \int (q!)^{-1}\, |z|^{2q}\, |u|^2\, dG \]
where $u \in S_* T \otimes V$ is thought of as a function with values in $V$ and $|v|^2 =$
$\sum v_i v_i^*$. Let $(\ ,\ )_q$ and $(\ ,\ )_{q*}$ be the associated hermitian forms. Let
$C_k{}^q$ be the binomial coefficients.

Lemma 2. If $u \in S_k T \otimes V$, then $\|u\|_q^2 = C_q{}^k\, \|u\|_0^2$.

Proof. By the binomial theorem
\[ \sum_{\gamma} C_{\gamma}{}^{\delta} X^{\gamma} = (1+X)^{\delta} \]
where $C_{\gamma}{}^{\delta} = \delta!\,(\gamma!\,(\delta-\gamma)!)^{-1}$ and $X=(X_1,\dots,X_n)$. Substituting $X_1 = \cdots =$
$X_n = t$ and comparing coefficients of $t^q$ we have
\[ \sum_{|\gamma|=q} C_{\gamma}{}^{\delta} = C_q{}^{|\delta|}. \tag{4} \]
Using (3) and (4) we have
\[ \begin{aligned} \|z^{\alpha} e_i\|_q^2 &= \sum_{|\gamma|=q} (\gamma!)^{-1}\, \|D^{\gamma} z^{\alpha}\|_0^2 = \sum_{|\gamma|=q} (\gamma!)^{-1}\, (\alpha!)^2\, ((\alpha-\gamma)!)^{-1} \\ &= \alpha!\, C_q{}^{|\alpha|} = C_q{}^k\, \|z^{\alpha} e_i\|_0^2 . \end{aligned} \]
Combining this with the readily verified fact that $z^{\alpha} e_i$ and $z^{\beta} e_j$ in $S_k T \otimes V$
are orthogonal for both $(\ ,\ )_q$ and $(\ ,\ )_0$ unless $\alpha=\beta$ and $i=j$, we obtain
the lemma.

Lemma 3. There is a constant $C$ depending only on $q$ such that
\[ \|u\|_q^2 \leqq \|u\|_{q*}^2 \leqq C \sum_{j \leqq q} \|u\|_j^2 \]
for every $u$ in $S_* T \otimes V$.

Proof. We shall need the commutation formulae
\[ \begin{aligned} (\alpha!\,\beta!)^{-1} z^{\alpha} D^{\beta} u &= \sum_{\gamma} (-1)^{|\gamma|}\, (\gamma!\,(\beta-\gamma)!\,(\alpha-\gamma)!)^{-1}\, D^{\beta-\gamma} z^{\alpha-\gamma} u \\ (\alpha!\,\beta!)^{-1} D^{\alpha} z^{\beta} u &= \sum_{\gamma} (\gamma!\,(\beta-\gamma)!\,(\alpha-\gamma)!)^{-1}\, z^{\beta-\gamma} D^{\alpha-\gamma} u \end{aligned} \tag{5} \]
where the first may be obtained as the coefficient of $s^{\alpha} t^{\beta}$ in the power
series expansion of the Weyl commutation relations
\[ e^{sz} e^{tD} u(z) = e^{-st} e^{tD} e^{sz} u(z) \]
\[ \text{i.e. } e^{sz} u(z+t) = e^{-st} e^{s(z+t)} u(z+t) \]

17*

\margin{52}

\clearpage
% Tapis de Quillen, with Quillen's offprint and preprint (cote 162-5, typed part): archive cote 162-5, site page 53 (PDF page 54).
% offprint. Draft compared: none.
% AI-transcribed by Claude Opus 5.5 (claude-opus-5-5, adaptive thinking, effort high), Message Batches API, batch t162_full (msgbatch_01UAh3zyWPd454AeDkhnmQ33), 2026-11-13; not yet checked by eye.
\scanpage{162-5/53}
\mspage{240} D.G. Quillen:

and where the second may be derived similarly (see also [2], p. 153). Thus
\[ \|u\|^2_{q*} = \sum_{|\alpha|=q} (\alpha!)^{-1}\, \|z^\alpha u\|^2_0 = \sum_{|\alpha|=q} (\alpha!)^{-1} (u, D^\alpha z^\alpha u)_0 \]
\[ = \sum_{|\alpha|=q} \sum_{\gamma} \alpha!\,(\gamma!\,(\alpha-\gamma)!^2)^{-1} (u, z^{\alpha-\gamma} D^{\alpha-\gamma} u)_0 \]
\[ = \sum_{|\alpha|=q} \sum_{\gamma} C_\gamma{}^\alpha ((\alpha-\gamma)!)^{-1}\, \|D^{\alpha-\gamma} u\|^2_0 \]
from which the lemma immediately follows.

Proof of the Theorem. The implication (ii)$\Rightarrow$(i) is clear, so assume (i)
holds. If $u \in S_{m+N}\, T \otimes V$ then by Lemma 1, (2), and (5) we have with
indices running like $1 \leqq i,\ j \leqq r$, $|\alpha| = |\beta| = m$
\[ E^N(u,u) = \sum a_{\alpha i, \beta j} (D^\alpha u_i, D^\beta u_j)_0 \]
\[ = \sum a_{\alpha i, \beta j} (u_i, z^\alpha D^\beta u_j)_0 \]
\[ = \int \sum (-1)^{|\gamma|}\, a_{\alpha i, \beta j}\, \alpha!\, \beta!\, (\gamma!\, (\beta-\gamma)!\, (\alpha-\gamma)!)^{-1}\, z^{\beta-\gamma} u_i (z^{\alpha-\gamma} u_j)^*\, dG \]
\[ = \int [E^\wedge (z^*, u) + Q(z,u)]\, dG \]
where $Q(z,u)$ is of degree $<m$ in $z_i z_j^*$ and of degree 1 in $u_i u_j^*$. Thus
for some constant $C$
\[ |Q(z,u)| \leqq C \sum_{i<m} (i!)^{-1}\, |z|^{2i}\, |u|^2, \]
while by hypothesis on $E$ there is an $\varepsilon > 0$ such that
\[ |E^\wedge (z^*, u)| \geqq \varepsilon (m!)^{-1}\, |z|^{2m}\, |u|^2. \]
Therefore
\[ E^N(u,u) \geqq \int (\varepsilon (m!)^{-1}\, |z|^{2m} - C \sum_{i<m} (i!)^{-1}\, |z|^{2i})\, |u|^2\, dG \]
\[ = \varepsilon\, \|u\|^2_{m*} - C \sum_{i<m} \|u\|^2_{i*} \]
\[ \geqq \varepsilon\, \|u\|^2_m - C' \sum_{i<m} \|u\|^2_i \]
using Lemma 3. By Lemma 2 we have $E^N(u,u) \geqq P(N)\, \|u\|^2_0$ where
\[ P(N) = \varepsilon\, C_m{}^{m+N} - C' \sum_{i<m} C_i{}^{i+N} \]
is a polynomial of degree $m$ in $N$ with positive leading coefficient. Thus
for $N$ sufficiently large $P(N) > 0$ so $E^N$ is positive definite and the theorem
is proved.

We shall now give some applications of the theorem when $V = \mathbf{C}$
\margin{\unsure{$\backslash$}} and $T = \mathbf{C}^n$.

\mspage{\hand{53}}

\clearpage
% Tapis de Quillen, with Quillen's offprint and preprint (cote 162-5, typed part): archive cote 162-5, site page 54 (PDF page 55).
% offprint. Draft compared: none.
% AI-transcribed by Claude Opus 5.5 (claude-opus-5-5, adaptive thinking, effort high), Message Batches API, batch t162_full (msgbatch_01UAh3zyWPd454AeDkhnmQ33), 2026-11-13; not yet checked by eye.
\scanpage{162-5/54}
\mspage{On the Representation of Hermitian Forms as Sums of Squares \quad 241}

Corollary 1. Let
\[ E^{\wedge}(\zeta)=\sum_{|\alpha|=|\beta|=m} a_{\alpha\beta}\,\zeta^{\alpha}\,\zeta^{*\beta} \]
be a hermitian form homogeneous of degree $m$ on $\mathbf{C}^n$ which is strictly
positive, i.e. $E^{\wedge}(\zeta)>0$ if $\zeta\neq 0$. Then for sufficiently large $N$,
\[ E^{\wedge}(\zeta)=\frac{\sum\limits_{i}|P_i(\zeta)|^2}{\sum\limits_{|\gamma|=N}(\gamma!)^{-1}\,|\zeta^{\gamma}|^2} \]
where $P_i(\zeta)$ is a basis for the holomorphic polynomials of degree $m+N$ on $\mathbf{C}^n$.

Corollary 2. Let
\[ E^{\sim}(f)=\sum_{|\alpha|=|\beta|=m} a_{\alpha\beta}\,D^{\alpha}f(D^{\beta}f)^* \]
be a differential quadratic form in $n$ variables homogeneous of degree
$(m,m)$, such that $a_{\beta\alpha}=a^*_{\alpha\beta}$. If $E^{\sim}$ has no non-zero characteristics, i.e.
$\sum a_{\alpha\beta}\,\zeta^{\alpha}\,\zeta^{*\beta}>0$ for $\zeta\neq 0$, then for sufficiently large $N$ we may write
\[ \sum_{|\gamma|=N}(\gamma!)^{-1}\,E^{\sim}(D^{\gamma}f)=\sum_{i}|P_i(D)\,f|^2 \]
where $P_i(D)$ is a basis for the homogeneous polynomials of degree $m+N$
in $D_1,\dots,D_n$.

These corollaries are just reinterpretations of the implication (i)$\Rightarrow$(ii)
of the theorem.

Example. Let $E^{\sim}(f)=|D_1^2\,u|^2-|D_1\,D_2\,u|^2+|D_2^2\,u|^2$ where $n=2$. The dif-%
ferential quadratic form $E^{\sim}$ is not a sum of squares because $E^{\sim}(z^1\,z^2)<0$.
But $E^{\sim}(\zeta)=(|\zeta_1|^2-|\zeta_2|^2)^2+|\zeta_1\,\zeta_2|^2>0$ if $\zeta\neq 0$ so the theorem applies.
In fact $E^{1\sim}(f)=|D_1^3\,f|^2+|D_2^3\,f|^2$ and $E^N$ is positive definite for $N\geqq 3$.

Corollary 3. (Projective Nullstellensatz over $\mathbf{C}$.) If $Q_j(z)$ $1\leqq j\leqq l$ are
homogeneous polynomials in $z^1,\dots,z^n$ without a common zero $\zeta\neq 0$, then
the ideal generated by the $Q_j$ contains all homogeneous polynomials of
sufficiently high degree.

Proof. We may assume that the $Q_j$ have the same degree $m$. Let
\[ E^{\wedge}(\zeta)=\sum_{j}|Q_j(\zeta)|^2; \]
then $E^{\wedge}(\zeta)>0$ if $\zeta\neq 0$ so by Corollary 2 for all sufficiently high $N$, there
is an expression
\[ \sum_{j}\sum_{|\gamma|=N}(\gamma!)^{-1}\,|D^{\gamma}\,Q_j(D)\,f|^2=\sum|P_i(D)\,f|^2 \]
where $P_i(D)$ is a basis for the space $W$ of homogeneous polynomials
of degree $m+N$. It follows that the set $D^{\gamma}\,Q_j(D)$ for $|\gamma|=N$, $1\leqq j\leqq l$
spans $W$ for otherwise we could take $f$ to be a non-zero polynomial in

\mspage{\hand{54}}

\clearpage
% Tapis de Quillen, with Quillen's offprint and preprint (cote 162-5, typed part): archive cote 162-5, site page 55 (PDF page 56).
% offprint. Draft compared: none.
% AI-transcribed by Claude Opus 5.5 (claude-opus-5-5, adaptive thinking, effort high), Message Batches API, batch t162_full (msgbatch_01UAh3zyWPd454AeDkhnmQ33), 2026-11-13; not yet checked by eye.
\scanpage{162-5/55}
\mspage{242} D.G. Quillen: On the Representation of Hermitian Forms as Sums of Squares

$z$ homogeneous of degree $m+N$ killed by the $D^\gamma Q_j(D)$ and get a contra-%
diction. This proves the corollary.

Corollary 3 implies the usual Nullstellensatz by means of the Rabino-%
witz trick and making polynomials homogeneous by adding an extra
variable. Hence we have given a direct analytical proof of the Null-%
stellensatz over $\mathbf{C}$ bearing some resemblance to the Liouville theorem
proof of the fundamental theorem of algebra.

Bibliography

1. Jacobson, N.: Lectures in abstract algebra III. Princeton: van Nostrand 1964.
2. Trèves, F.: Linear partial differential equations with constant coefficients. New York:
Gordon \& Breach 1966.

Dr. D. G. Quillen
Massachusetts Institute of Technology
Department of Mathematics
Cambridge, Mass. 02139, USA

(Received February 27, 1968)

\mspage{\hand{55}}

\clearpage
% Tapis de Quillen, with Quillen's offprint and preprint (cote 162-5, typed part): archive cote 162-5, site page 57 (PDF page 58).
% preprint. Draft compared: none.
% AI-transcribed by Claude Opus 5.5 (claude-opus-5-5, adaptive thinking, effort high), Message Batches API, batch t162_full (msgbatch_01UAh3zyWPd454AeDkhnmQ33), 2026-11-13; not yet checked by eye.
\scanpage{162-5/57}
% NOTE: at the top left there is a short typed or ruled dash and some ink specks, with no readable content.
\begin{center}
\uline{On a conjecture of Adams}*

\hand{D.G. Quillen}
\end{center}

This conjecture asserts that if $X$ is a finite CW complex
and $x \in K(X)$ , then for $n$ sufficiently large $p^{n}(\hand{$\psi^{p}$}x - x)$ repre-%
sents $0$ in $J(X)$ . Here $p$ is a prime number and $K$ is complex
$K$ theory. (Actually Adams states the conjecture when $p$ is any
integer and for real $K$ theory. It is clearly sufficient to prove
the result when \hand{$p$} is prime. We shall restrict ourselves entirely
to complex vector bundles for our methods do not apply to real $K$
theory.) Adams \hand{[1]} proves this conjecture if $x$ is a virtual sum
of line bundles \hand{and} reduces the conjecture for some other spaces
such as spheres \hand{to} this case. In this note we show how the etale
homotopy theory recently developed by Artin and Masur [2], which
extends the \unsure{etale} cohomology theory of Artin and Grothendieck [3],
may in principle be used to give a proof of this conjecture. Actu-%
ally there are certain results for the etale homotopy theory which
we need and which at the moment have not been proved, so this note
is really a sketch of a potential proof of this conjecture. How-%
ever the possibility of proving a theorem about CW complexes
using algebraic varieties in characteristic $p$ seems sufficiently
startling to merit publication in this incomplete form.

We begin by describing our key idea. In [1] Adams shows that
two vector bundles \hand{$E$} and \hand{$F$} over $X$ represent the same element
of $J(X)$ mod $p$-torsion if there exists a fiber preserving map

\noindent\rule{0.4\textwidth}{0.4pt} \margin{\unsure{NSF GP 6166}}

*This research \hand{was} supported \xout{\unsure{in part}} by \xout{\illegible{3}}
\xout{\illegible{4}}.

\mspage{57}

\clearpage
% Tapis de Quillen, with Quillen's offprint and preprint (cote 162-5, typed part): archive cote 162-5, site page 58 (PDF page 59).
% preprint. Draft compared: none.
% AI-transcribed by Claude Opus 5.5 (claude-opus-5-5, adaptive thinking, effort high), Message Batches API, batch t162_full (msgbatch_01UAh3zyWPd454AeDkhnmQ33), 2026-11-13; not yet checked by eye.
\scanpage{162-5/58}
$f:E - 0 \rightarrow F - 0$ which is of degree a power of $p$ on each fiber.
When $E$ is a line bundle, then $\hand{\psi}^pE = E^{\otimes p}$ and $f(v) = v^{\otimes p}$ is
the desired map to prove the conjecture for $E$ . Unfortunately
when $E$ has dimension greater than one, then $\hand{\psi}^pE$ is only a vir-%
tual bundle so the method fails. However if $E$ is an algebraic
vector bundle (locally free sheaf of finite rank) over a prescheme
$V$ of characteristic $p$ , then in virtue of the Frobenius homomor-%
phism there is an algebraic vector bundle $E^{(p)}$ whose class in
$K_{alg}(V)$ represents $\hand{\psi}^pE$ .\hand{*} If $E$ is given by a cocycle $(g_{ij})$ of
$n\mathrm{x}n$ matrices with respect to an open covering of $V$ , then $E^{(p)}$
is given by the cocycle $(g_{ij}^{(p)})$ , where $g_{ij}^{(p)}$ denotes the matrix
whose entries are the p-th powers of the corresponding entries of
$g_{ij}$ . Moreover there is a canonical map $F:E \rightarrow E^{(p)}$ given in lo-%
cal coordinates by $(X_i) \rightarrow (X_i^p)$ which is of (inseparable) degree
$p^{\hand{\dim E}}$ on each fiber. Hence the hypotheses for Adams' criterion is
verified in the algebraic situation and it remains to reduce the
topological situation to the algebraic one.

It evidently suffices to prove the conjecture when $X$ is the
Grassmannian of complex $N$ planes in $n + N$ space and where
$x \in K(X)$ is the class of the canonical quotient bundle. Let $G_Z$
be the Grassmannian scheme over Spec Z representing the functor
``isomorphism classes of n-dimensional vector bundles with $n + N$
generating sections'' and let $Q_Z$ be the canonical n-dimensional
bundle. Then $X = G_{C,cl}$ and $x$ is the class of $Q_{C,cl}$ , where

\hand{* That $E^{(p)}$ represents $\psi^pE$ follows since $E \mapsto E^{(p)}$ defines a
homomorphism $K_{alg} \rightarrow K_{alg}$ which coincides with $\psi^p$ on line bundles,
together with the splitting principle: $K_{alg}(PE) \unsure{\cong|=} K_{alg}(X)[T]/(\lambda_{-T}(E))$.}

\mspage{58}

\clearpage
% Tapis de Quillen, with Quillen's offprint and preprint (cote 162-5, typed part): archive cote 162-5, site page 59 (PDF page 60).
% preprint. Draft compared: none.
% AI-transcribed by Claude Opus 5.5 (claude-opus-5-5, adaptive thinking, effort high), Message Batches API, batch t162_full (msgbatch_01UAh3zyWPd454AeDkhnmQ33), 2026-11-13; not yet checked by eye.
\scanpage{162-5/59}
\mspage{- 3 -}

$G_A,Q_A$ denote the base extension of $G_Z,Q_Z$ by the map Spec A $\to$
Spec Z , and where if Y is a prescheme of finite type over Spec C,
$Y_{cl}$ denotes the set of complex points of Y with its ``classical''
topology. Let R be a strict localization of Z at the prime p
and choose an embedding R $\to$ C . Then R is a henselian discrete
valuation ring whose residue field k is algebraically closed of
characteristic p . Let $\underline{H}$ be the homotopy category of connected
pointed CW complexes, and let $X \to \hat{X}$ be the \uline{etale homotopy type}
functor [2] from the category of geometrically pointed locally con-%
nected preschemes to pro-$\underline{H}$. Let C be the class of finite groups
of order prime to p , and let $\underline{X} \to \hat{\underline{X}}$ be the \uline{C-completion} functor
[2] pro $\underline{H} \to$ pro C$\underline{H}$ . Choose a rational point of $G_Z$ over Spec Z
and rationally point $G_A$ by base extension. Then as $G_R$ is proper
and smooth over Spec R and 1-connected, the fundamental comparison
theorems of [2] imply that there are isomorphisms in pro-C$\underline{H}$

\[
(1) \qquad G^{\wedge}_{C,cl} \xrightarrow[\sim]{u} G^{\wedge}_{C,et} \xleftarrow[\sim]{i} G^{\wedge}_{R,1,et} \overset{v}{\simeq} G^{\wedge}_{R,2,et} \xleftarrow[\sim]{j} G^{\wedge}_{\unsure{R|k},et}
\]

where i and j are base-change morphisms, u comes from the cano-%
nical map from the classical to the etale topology, \unsure{where|wiere} $G_{R,1}$
(resp. $G_{R,2}$) is the scheme $G_R$ with geometric point Spec C $\to$
Spec R $\to G_R$ (resp. Spec k $\to$ Spec R $\to G_R$) , and where v is the
isomorphism coming from shifting basepoints from 1 to 2 along the
``arc'' Spec R $\to G_R$ .

\mspage{\hand{59}}

\clearpage
% Tapis de Quillen, with Quillen's offprint and preprint (cote 162-5, typed part): archive cote 162-5, site page 60 (PDF page 61).
% preprint. Draft compared: none.
% AI-transcribed by Claude Opus 5.5 (claude-opus-5-5, adaptive thinking, effort high), Message Batches API, batch t162_full (msgbatch_01UAh3zyWPd454AeDkhnmQ33), 2026-11-13; not yet checked by eye.
\scanpage{162-5/60}
\mspage{- 4 -}

We shall now \uline{assume} \deleted{for the moment} the possibility of defining
for each object $\underline{X}$ of pro-$\underline{\mathrm{CH}}$ an abelian monoid $SF(\underline{X})$ of stable
isomorphism classes of oriented sphere fibrations with base $\underline{X}$ ,
such that there is a commutative diagram of monoids

(2)
\begin{diagramblock}
\begin{tikzcd}[column sep=small]
K(G_{C,cl}) \arrow[d, "J"'] & & & K_{alg}(G_R) \arrow[lll, "aj^!"'] \arrow[rr, "i^!"] \arrow[dl, "J"'] \arrow[dr, "J"] & & K_{alg}(G_k) \arrow[d, "J"] \\
SF(G_{C,cl}) \arrow[r, "b"] & SF(\hat{G}_{C,cl}) & SF(\hat{G}_{R,1,et}) \arrow[l, "u^*j^*"'] \arrow[rr, phantom, "\overset{v^*}{\cong}"] & & SF(\hat{G}_{\unsure{R|K},2,et}) \arrow[r, "i^*"] & SF(\hat{G}_{k,et})
\end{tikzcd}
\end{diagramblock}

where $a:K_{alg}(G_C) \to K(G_{C,cl})$ is induced by the map associating to
an algebraic vector bundle $E$ the topological bundle $E_{cl}$ , where
$b$ is induced by sending $S^q \to W \to G_{C,cl}$ into $\hat{S}^q \to \hat{W} \to \hat{G}_{C,cl}$ ,
and where $J$ is the map induced by sending an algebraic vector
bundle $E$ over $X$ into the sphere fibration $(E - 0)_{et} \to X_{et}$
which we denote by \unsure{$\$E$} .

With these assumptions we may now prove Adams' conjecture.
Denote by $cl(E)$ the class in $K$ or $K_{alg}$ of a vector bundle $E$ .
By the lemma at the end there is a vector bundle $E_R$ over $G_R$ with

\[ (3) \qquad cl(E_R) - n = \psi^p(clQ_R) - cl(Q_R) \]

Applying $i^!$ we obtain $cl(E_k) - n = \psi^p(clQ_k) - cl(Q_k) =$

\mspage{\hand{60}}

\clearpage
% Tapis de Quillen, with Quillen's offprint and preprint (cote 162-5, typed part): archive cote 162-5, site page 61 (PDF page 62).
% preprint. Draft compared: none.
% AI-transcribed by Claude Opus 5.5 (claude-opus-5-5, adaptive thinking, effort high), Message Batches API, batch t162_full (msgbatch_01UAh3zyWPd454AeDkhnmQ33), 2026-11-13; not yet checked by eye.
\scanpage{162-5/61}
\mspage{- 5 -}

$\mathrm{cl}(Q_k^{(p)}) - \mathrm{cl}(Q_k)$ .  But  $F:Q_k - 0 \to Q_k^{(p)} - 0$  is fiber-preserving
and it is purely inseparable, hence a homeomorphism for the etale
topology.  Hence  $\$Q_k^{(p)} \simeq \$Q_k$  so  $i^*J(\mathrm{cl}E_R) = J(\mathrm{cl}(E_k) - n) = 0$ .
As  $i^*$ is an isomorphism by (1), $J(\mathrm{cl}E_R) = 0$, so from the diagram
(2) we find  $bJ\mathrm{cl}(E_{C,\mathrm{cl}}) = 0$ .  But this means that
\[ S^{2n-1} \to (E_{C,\mathrm{cl}} - 0)^{\wedge} \to G^{\wedge}_{C,\mathrm{cl}} \qquad , \]
is a stably trivial fibration, that is, there is a map  $u$  in
$\{(E_{C,\mathrm{cl}} - 0)^{\wedge} , S^{2n-1}\}$  such that  $u\hat{\epsilon} = \mathrm{id}$ .  Here  $\{ \; , \; \}$  denotes
S-maps in pro-$\underline{CH}$  [2] \unsure{§8|§§}, and  $\epsilon:S^{2n-1} \to E_{C,\mathrm{cl}} - 0$  is the injection
of a fiber.  Consider the diagram

\begin{diagramblock}
\begin{tikzcd}[column sep=small]
\{(E_{C,\mathrm{cl}}-0) , S^{2n-1}\}^{\wedge} \arrow[r, "(\epsilon^*)^{\wedge}"] \arrow[d] & \{S^{2n-1} , S^{2n-1}\}^{\wedge} \arrow[r, "v"] \arrow[d, "\hand{}"] & \mathbb{Z}/r\mathbb{Z} \arrow[r] \arrow[d] & 0 \\
\{(E_{C,\mathrm{cl}}-0)^{\wedge}, \widehat{S^{2n-1}}\} \arrow[r, "(\hat{\epsilon})^*"] & \{\widehat{S^{2n-1}} , \widehat{S^{2n-1}}\} \arrow[r] & 0 &
\end{tikzcd}
\end{diagramblock}
% NOTE: the middle vertical arrow (from {S^{2n-1}, S^{2n-1}}^∧ down to {Ŝ^{2n-1}, Ŝ^{2n-1}}) is drawn by hand in darker ink: \hand{$\downarrow$}

where  $v$  is the cokernel of  $\epsilon^*$ , and where  $\wedge$  in the first row
denotes the pro-C completion of an abelian group, which here is  $\otimes\hat{Z}$ ,
since the abelian groups are finitely generated, and hence the first
row is exact.  By [2] \unsure{§8|§§}, the first two vertical maps are isomorphisms,
hence  $(\mathbb{Z}/r\mathbb{Z})^{\wedge} = 0$  and  $r = p^{\unsure{i}}$ .  This means that there is an S-map
$w:(E_{C,\mathrm{cl}} - 0) \to S^{2n-1}$  such that  $w\epsilon$  is of degree $p^{\unsure{i}}$ , hence for
sufficiently large  N  Adams' criterion may be applied to the bundles

\mspage{\hand{61}}

\clearpage
% Tapis de Quillen, with Quillen's offprint and preprint (cote 162-5, typed part): archive cote 162-5, site page 62 (PDF page 63).
% preprint. Draft compared: none.
% AI-transcribed by Claude Opus 5.5 (claude-opus-5-5, adaptive thinking, effort high), Message Batches API, batch t162_full (msgbatch_01UAh3zyWPd454AeDkhnmQ33), 2026-11-13; not yet checked by eye.
\scanpage{162-5/62}
\mspage{- 6 -}

$E_{C,cl} \oplus N$ and $n \oplus N$ . Hence there is a $k$ such that $0 =$
$p^{k}\, J(clE_{C,cl}) = p^{k}\, J(\psi^{p}x - x)$ where $x = clQ_{C,cl}$ and the con-%
jecture is proved.

We shall now discuss the assumptions used in the above proof,
proving what we can. There are three problems:

A. Define oriented sphere fibrations over objects of pro-$\underline{CH}$ .

B. Show that if $E$ is a vector bundle over the pointed pre-%
scheme $X$ , where $X = G_{k}$ , $G_{R,1}$ or $G_{R,2}$ , then $E - 0 \to X$ deter-%
mines an oriented sphere fibration $\unsure{\$}E$ over $X^{\wedge}_{et}$ .

C. Show that the mapping $\unsure{\$}{:}\mathrm{Vect}\,X \to SF(\overset{\wedge}{X}_{et})$ induced a map
$J{:}K_{alg}(X) \to SF(X^{\wedge}_{et})$ , where Vect $X$ denotes the set of isomorphism
classes of vector bundles over $X$ .

We now give what we believe to be a solution of $A$ , we give
evidence for $B$ , and finally we will solve $C$.

$\underline{A}$: Let $\underline{X}$ be an object of pro-$\underline{CH}$ . Define an oriented
$(n - 1)$-sphere fibration with base $\underline{X}$ to be a triple $(\underline{f},u,v)$
where $\underline{f} = \{f_{i}{:}E_{i} \to B_{i}\ ,\ i \in I\}$ is a pro-object in the homotopy
category of maps in $\underline{CH}$ , where $u{:}\{B_{i}\} \to \underline{X}$ is an isomorphism,
and where $v{:}(S^{n-1})^{b} \xrightarrow{\sim} \{Ef_{i}\}^{b}$ . Here $Ef = E\times_{B}B^{I}\times_{B}(*)$ is the
fiber of the Hurewicz fibration associated to a map $f$ , and ``b''
denotes the ``flat'' functor of [2], §3. Then the set $SF_{n}(\underline{X})$ of
isomorphism classes of $(n - 1)$-sphere fibrations with base $\underline{X}$
is a contravariant functor of $\underline{X}$ and the fiber-join [4] operation
defines a map $SF_{n}(\underline{X})\times SF_{m}(\underline{X}) \to SF_{n+m}(\underline{X})$ , permitting us to define

\mspage{\hand{62}}

\clearpage
% Tapis de Quillen, with Quillen's offprint and preprint (cote 162-5, typed part): archive cote 162-5, site page 63 (PDF page 64).
% preprint. Draft compared: none.
% AI-transcribed by Claude Opus 5.5 (claude-opus-5-5, adaptive thinking, effort high), Message Batches API, batch t162_full (msgbatch_01UAh3zyWPd454AeDkhnmQ33), 2026-11-13; not yet checked by eye.
\scanpage{162-5/63}
\mspage{- 7 -}

the abelian monoid of stable oriented sphere fibrations with base
$\underline{X}$ in the usual way. By [2], th. 4.5, C-completion defines a
natural transformation $[X,BSF_n] \to SF_n(\hat{X})$ if $X$ is a 1-connected
object of $\underline{H}$ , where $BSF_n$ classifies oriented $(n - 1)$-sphere
fibrations in $\underline{H}$ . We conjecture that $[\underline{X},\widehat{BSF}_n] \simeq SF_n(\underline{X})$ but we
will not need this fact.

\underline{B}: By an easy spectral sequence argument one sees that up to
isomorphism an oriented $(n - 1)$-sphere fibration $(\underline{f},u,v)$ over $\underline{X}$
may be identified with the triple $(\underline{f},u,U)$ where $U \in H^n(C\underline{f},\hat{Z})$ is
the Thom class of the sphere fibration. Here $C\underline{f} = \{Cf_i, i \in I\}$ ,
$Cf_i$ is the cone on the map $f_i$ , and for $U$ to be the Thom class
means that $\cup U : H^q(X,\hat{Z}) \simeq H^{q+n}(C\underline{f},\hat{Z})$ and that $(\tilde{v})^*U$ is the cano-%
nical generator of $H^n(\hat{S}^n,\hat{Z})$ where $\tilde{v} : \hat{S}^n \to C\underline{f}$ is the map induced
by $v$ .

Now let $E$ be a vector bundle over $X$ where $X = G_R$ or $G_k$ .
Then there is an exact sequence

\[ (4) \qquad \to H^q_X(E,\hat{Z}) \to H^q(E,\hat{Z}) \to H^q(E - 0,\hat{Z}) \unsure{\hat{\Delta}} \]

where $H^q_X(E,\hat{Z})$ is the local cohomology of $E$ along the 0 section.
By [3], Expose XVI, remark (3.10), there is a canonical class
$c_n \in H^{2d}_X(E,\mu_n^{\otimes d})$ , where $(n,p) = 1$ and $d = \dim E$ , which gives
rise to an isomorphism

\margin{63}

\clearpage
% Tapis de Quillen, with Quillen's offprint and preprint (cote 162-5, typed part): archive cote 162-5, site page 64 (PDF page 65).
% preprint. Draft compared: none.
% AI-transcribed by Claude Opus 5.5 (claude-opus-5-5, adaptive thinking, effort high), Message Batches API, batch t162_full (msgbatch_01UAh3zyWPd454AeDkhnmQ33), 2026-11-13; not yet checked by eye.
\scanpage{162-5/64}
\mspage{- 8 -}

\margin{\checkmark}\[ UC_n : H^q(X, Z_n) \to H^{q+2d}_X (E, \mu_n^{\otimes d}) . \]

When $X^\prime = \unsure{G_k|O_k}$ , the $\mu_n$ in this formula refers to the abelian
group $\mu_{n,k}$ of nth roots of unity in $k$ . But the map $k^* \unsure{\overset{s}{\leftarrow}}$
$R^* \to C^*$ where $s$ is the unique multiplicative section gives a
canonical isomorphism $\mu_{n,k} \simeq \mu_{n,C} \simeq Z_n$ where the last isomorphism
comes from the distinguished nth root exp$(2\pi i/n)$. When $X = \unsure{G_R|O_R}$
then $\mu_{n,R} = \mu_{n,C} \simeq Z_n$ . Hence there is a canonical element
$C = \{C_n\} \in H^{2d}_X(E, \hat{Z})$ giving an isomorphism

(5) \[ UC : H^q(X, \hat{Z}) \simeq H^{q+2d}_X(E, \hat{Z}) . \]

Now choose a basepoint of $E - 0$ lying over the basepoint of $X$ .
I have been told that by careful analysis of the etale homotopy
type functor it is possible to rigidify $(E - 0)^\wedge_{et} \to E^\wedge_{et}$ to be-%
come a pro-object $\underline{f}$ in the homotopy category of maps in $\underline{CH}$ .
Let us assume that further analysis will show there is a canonical
isomorphism $H^*(C\underline{f}, \hat{Z}) \simeq H^*_X(E, \hat{Z})$ which is reasonable in virtue of
(4). Then (5) gives a Thom class $U \in H^{2d}(C\underline{f}, \hat{Z})$ showing that
$(\underline{f},u,U)$ is an oriented sphere fibration \$E with base $X^\wedge_{et}$ ,
where $u : \hat{X}_{et} \simeq \hat{E}_{et}$ is the isomorphism coming from the projection
$E \to X$ . As (5) is compatible with base-change $g^*(\$E) \simeq \$(g^*E)$
when $g$ is a map of preschemes.

\margin{64}

\clearpage
% Tapis de Quillen, with Quillen's offprint and preprint (cote 162-5, typed part): archive cote 162-5, site page 65 (PDF page 66).
% preprint. Draft compared: none.
% AI-transcribed by Claude Opus 5.5 (claude-opus-5-5, adaptive thinking, effort high), Message Batches API, batch t162_full (msgbatch_01UAh3zyWPd454AeDkhnmQ33), 2026-11-13; not yet checked by eye.
\scanpage{162-5/65}
\mspage{- 9 -}

\uline{C}: It is necessary to show that the image of \$:Vect X $\rightarrow$
SF($\hat{X}_{et}$) lies in the subgroup of nivertible elements of the mo-%
noid SF($\hat{X}_{et}$) and that if

\[ (6) \qquad 0 \rightarrow E' \rightarrow E \rightarrow E'' \rightarrow 0 \]

is an exact sequence of vector bundles over $X$ , then \$E =
\$E' + \$E" . If $X = G_R$ with basepoint 1 or 2, this is clear from
the commutative diagram

\begin{diagramblock}\begin{tikzcd}[column sep=small]
K(G_{C,cl}) \arrow[d, "J"'] & & K_{alg}(G_R) \arrow[l, "\unsure{\leftarrow|\rightarrow}" description, phantom] \arrow[dl] \arrow[d] & Vect(G_R) \arrow[l] \arrow[d, "\$"'] \arrow[dr, "\$"] & \\
SF(G_{C,cl}) \arrow[r] & SF(\hat{X}_{C,cl}) \arrow[r, phantom, "\cong"] & SF(\hat{G}_{R,1,et}) \arrow[rr, phantom, "\overset{\Sigma^*}{\cong}"] & & SF(\hat{G}_{R,2,et})
\end{tikzcd}\end{diagramblock}
\margin{\unsure{$\xi$ Win}}

We cannot handle $G_k$ in this way because $i^*$:Vect($G_R$) $\rightarrow$ Vect($G_k$)
is not surjective. But consider the following trick. Suppose
that (6) is an exact sequence of vector bundles on $G_k$ and that
E is generated by sections. Let $D_Z = D_Z^{N,p,q}$ be the drapeaux
scheme over Spec Z which classifies diagrams

\[ N \rightarrow E_1^p \rightarrow E_2^q \]

of vector bundles over any prescheme, where $p$ = dim E , $q$ = dim E"

\mspage{\hand{65}}

\clearpage
% Tapis de Quillen, with Quillen's offprint and preprint (cote 162-5, typed part): archive cote 162-5, site page 66 (PDF page 67).
% preprint. Draft compared: none.
% AI-transcribed by Claude Opus 5.5 (claude-opus-5-5, adaptive thinking, effort high), Message Batches API, batch t162_full (msgbatch_01UAh3zyWPd454AeDkhnmQ33), 2026-11-13; not yet checked by eye.
\scanpage{162-5/66}
\mspage{- 10 -}

and $N$ is the number of sections required to generate $E$ . On
$D_Z$ there is a canonical exact sequence

\[ (7) \qquad 0 \to Q'_Z \to Q_Z \to Q''_Z \to 0 \]

and there is a map $f:G_k \to D_k$ which induces the exact sequence
(6) from the base extension of (7) to $k$ . As $D_R$ is smooth and
proper over Spec $R$ and 1-connected all arguments used above for
$G_R$ apply, that is $\$:\mathrm{Vect}\ D_R \to \hat{\mathrm{SF}}(\hat{D}_{R,2,et})$ induces a map $J$ .
Hence since

\begin{diagramblock}
\begin{tikzcd}
\mathrm{Vect}\ D_R \arrow[r, "\unsure{i^*|1^*}"] \arrow[d, "\$"] & \mathrm{Vect}\ D_k \arrow[r, "f^*"] \arrow[d, "\$"] & \mathrm{Vect}\ G_k \arrow[d, "\$"] \\
\mathrm{SF}(\hat{D}_{R,2,et}) \arrow[r, "\unsure{i^*|1^*}"] & \mathrm{SF}(\hat{D}_{k,et}) \arrow[r, "f^*"] & \mathrm{SF}(\hat{G}_{k,et})
\end{tikzcd}
\end{diagramblock}

Therefore $\$$ is additive for exact sequences (6) over $G_k$
when $E$ is generated by its sections. Call such a bundle \uline{full}.
If $E$ is full then there is an exact sequence $0 \to F \to n \to E \to 0$ ,
so $\$E + \$F = \$n = 0$ and $\$E$ is an invertible element of
$\mathrm{SF}(\hat{G}_{k,et})$ . If $E$ is any vector bundle on $G_k$ , then expressing
$E^*$ as a quotient of $mL^{-k}$ , where $L$ is the line bundle of the
canonical projective embedding, we obtain an exact sequence $0 \to$
$E \to F_1 \to F_2 \to 0$ with $F_1$ full. Hence $\$E + \$F_2 = \$F_1$ and
$\$E + (SF_2 - SF_1) = 0$ showing that $\$E$ is invertible for any $E$

\margin{66}

\clearpage
% Tapis de Quillen, with Quillen's offprint and preprint (cote 162-5, typed part): archive cote 162-5, site page 67 (PDF page 68).
% preprint. Draft compared: none.
% AI-transcribed by Claude Opus 5.5 (claude-opus-5-5, adaptive thinking, effort high), Message Batches API, batch t162_full (msgbatch_01UAh3zyWPd454AeDkhnmQ33), 2026-11-13; not yet checked by eye.
\scanpage{162-5/67}
\mspage{- 11 -}

in Vect $G_k$ . Finally given any exact sequence \unsure{(6),it|(6);it} may be
embedded in a diagram

\begin{diagramblock}
\begin{tikzcd}
 & & 0 & 0 & \\
 & & F_3 \arrow[u] \arrow[r, equal] & F_3 \arrow[u] & \\
0 \arrow[r] & E' \arrow[r] & F_1 \arrow[u] \arrow[r] & F_2 \arrow[u] \arrow[r] & 0 \\
0 \arrow[r] & E' \arrow[u, equal] \arrow[r] & E \arrow[u] \arrow[r] & E'' \arrow[u] \arrow[r] & 0 \\
 & & 0 \arrow[u] & 0 \arrow[u] &
\end{tikzcd}
\end{diagramblock}

with $F_1$ , hence $F_2$ full. Then $\$E = \$F_1 - \$F_3 = (\$E' + \$F_2) -$
$(\$F_2 - \$E'') = \$E' + \$E''$ which completes the proof of C.
It remains to prove the following result needed for (3).

\uline{Lemma}: Let $X$ be a noetherian prescheme \unsure{hawing|having} a very
ample line bundle $L$ . Then any $x \in K_{alg}(X)$ may be expressed
$x = cl(E) - n$ where $E$ is a vector bundle over $X$ .

\uline{Proof}: $K_{alg}(X)$ is the free abelian group generated by Vect $X$
with relations $cl(E) = cl(E') + cl(E'')$ for any exact sequence (6).
As $cl(E \oplus F) = cl(E) + cl(F)$ one may write $x = cl(E_1) - cl(E_2)$ .
As $L$ is ample there are exact sequences
\[ 0 \to E_3 \to n \to L^k \otimes \unsure{E_2^*} \to 0 \]
\[ 0 \to \unsure{E_4} \to m \to L^k \to 0 \]

\margin{67}

\clearpage
% Tapis de Quillen, with Quillen's offprint and preprint (cote 162-5, typed part): archive cote 162-5, site page 68 (PDF page 69).
% preprint. Draft compared: none.
% AI-transcribed by Claude Opus 5.5 (claude-opus-5-5, adaptive thinking, effort high), Message Batches API, batch t162_full (msgbatch_01UAh3zyWPd454AeDkhnmQ33), 2026-11-13; not yet checked by eye.
\scanpage{162-5/68}
\mspage{- 12 -}

for suitable $k,m,n > 0$ . Hence $cl(E_2) = n\,cl(L^k) - cl(L^k \unsure{\otimes|\oplus} E_3^{\unsure{*}})$
and $x = cl(E_1) + cl(L^k \unsure{\otimes|\oplus} E_3^{\unsure{*}}) - n(m - cl(E_4)) = cl(E_1 \unsure{\oplus|\otimes} L^k \unsure{\otimes|\oplus} E_3^{\unsure{*}} \unsure{\oplus|\otimes} nE_4)$
$- nm$ .
% NOTE: a smudged typed dashed line runs across the page below this line; no characters can be read in it

\uline{Bibliography}

[1] Adams, J. F., On the groups J(X)-I, Topology Vol. 2 (1963)
p.p. 181-195.

[2] Artin, M. and Mazur, B., Homotopy of Varieties in the
Etale Topology (to appear).

[3] Artin, M. and Grothendieck, A., S\hand{\'e}minare de G\hand{\'e}om\hand{\'e}trie
algebrique (63-64) - Cohomologie \hand{\'e}tale des Schemas.
(Mimeographed notes, Institute des Hautes \hand{\'E}tudes \hand{'}
Scientifiques, 1963-64).

[4] Hall, I. M., The Generalized Whitney Sum, Quart. J. Math.
Oxford (2), 16 (1965), 360-84.

Massachusetts Institute of Technology

\margin{68}

\end{document}
